% ===========================================================================
%  Sampling Geometries and Band-Limits for Radiance Fields
%  B.Tech Project · Department of Electrical Engineering · IIT Kharagpur
%
%  Build:  make thesis      (from the repository root)
%  which regenerates thesis/generated/*.tex from results/runs.csv first, so
%  no measurement in this document is ever typed by hand.
% ===========================================================================
\documentclass[11pt,a4paper,oneside]{report}
% ===========================================================================
%  Preamble. No measurement appears in this file or in any chapter file —
%  every measured number is \input from generated/measured.tex, which
%  tools/make_tex.py writes from results/runs.csv.
% ===========================================================================
\usepackage[T1]{fontenc}
\usepackage[utf8]{inputenc}
\usepackage{lmodern}
\usepackage{microtype}

\usepackage[a4paper,margin=27mm,bottom=30mm]{geometry}
\usepackage{amsmath,amssymb,amsthm,bm}
\usepackage{graphicx}
\usepackage{booktabs}
\usepackage{array}
\usepackage{multirow}
\usepackage{caption}
\usepackage{enumitem}
\usepackage{xcolor}
\usepackage{fancyhdr}
\usepackage{titlesec}
\usepackage[hidelinks,breaklinks]{hyperref}

\definecolor{ink}{HTML}{1A1D21}
\definecolor{ink2}{HTML}{52514E}
\definecolor{accent}{HTML}{2A78D6}
\definecolor{warn}{HTML}{EB6834}
\definecolor{rule}{HTML}{D8DDE3}

\hypersetup{colorlinks=true, linkcolor=ink, citecolor=accent, urlcolor=accent}

\captionsetup{font=small, labelfont={bf,color=ink}, textfont={color=ink2},
              width=0.94\linewidth, skip=6pt}
\captionsetup[table]{position=top, skip=4pt}

\setlength{\parskip}{0.45em}
\setlength{\parindent}{0pt}
\renewcommand{\arraystretch}{1.15}

\pagestyle{fancy}
\fancyhf{}
\renewcommand{\headrulewidth}{0.4pt}
\fancyhead[L]{\footnotesize\color{ink2} Sampling-Geometry Band-Limits for Radiance Fields}
\fancyhead[R]{\footnotesize\color{ink2} Rickarya Das \textperiodcentered\ 23EE30019}
\fancyfoot[C]{\thepage}

\titleformat{\chapter}[display]
  {\normalfont\Large\bfseries\color{ink}}
  {\color{accent}\normalsize\MakeUppercase{\chaptertitlename}\ \thechapter}{6pt}{\LARGE}
\titlespacing*{\chapter}{0pt}{0pt}{18pt}

% --------------------------------------------------------------- semantics
% A claim the project has measured, versus one it has only predicted. Keeping
% these visually distinct is the whole point of §7's discipline.
\newcommand{\pass}{\textcolor{accent}{\textbf{pass}}}
\newcommand{\fail}{\textcolor{warn}{\textbf{FAIL}}}
\newcommand{\skipped}{\textcolor{ink2}{skip}}

\newenvironment{claimbox}[1]
  {\par\vspace{4pt}\noindent\begin{minipage}{\linewidth}%
   \setlength{\fboxsep}{9pt}%
   \colorlet{shade}{accent!6}%
   \begin{tcbnone}%
   \noindent\textcolor{accent}{\footnotesize\textbf{\MakeUppercase{#1}}}\par\vspace{2pt}%
   \itshape}
  {\end{tcbnone}\end{minipage}\par\vspace{4pt}}
\newenvironment{tcbnone}{}{}

% A framed statement, used for the thesis claim and the two propositions.
\newcommand{\keybox}[2]{%
  \par\vspace{6pt}%
  {\color{rule}\hrule height 0.8pt}%
  \vspace{5pt}%
  \noindent\textcolor{accent}{\footnotesize\textbf{\MakeUppercase{#1}}}\par\vspace{3pt}%
  \noindent #2\par
  \vspace{5pt}{\color{rule}\hrule height 0.8pt}\par\vspace{6pt}}

\newcommand{\status}[1]{\textcolor{warn}{\textbf{#1}}}

\newtheorem{proposition}{Proposition}


% Every measured number in this document. Written by tools/make_tex.py from
% results/runs.csv; an unmeasured quantity expands to a visible marker, never to
% a plausible-looking number.
\input{generated/measured}

\begin{document}

% --------------------------------------------------------------- title page
\thispagestyle{empty}
\begin{center}
  \vspace*{12mm}
  {\footnotesize\color{ink2}\MakeUppercase{B.Tech Project \textperiodcentered\ Autumn 2026 -- Spring 2027}}\\[2pt]
  {\footnotesize\color{ink2}\MakeUppercase{Department of Electrical Engineering \textperiodcentered\ IIT Kharagpur}}

  \vspace{22mm}
  {\Huge\bfseries\color{ink} Sampling Geometries\\[4pt] and Band-Limits for\\[4pt] Radiance Fields\par}

  \vspace{10mm}
  \begin{minipage}{0.82\linewidth}\centering\itshape\color{ink2}
    The size limit on a splatting primitive should be the inverse of the
    information the capture accumulated about it, floored at the per-view
    Nyquist limit --- a matrix, not a scalar.
  \end{minipage}

  \vspace{20mm}
  \begin{tabular}{r@{\hspace{8pt}}l}
    \color{ink2}\footnotesize Candidate & Rickarya Das \textperiodcentered\ 23EE30019\\
    \color{ink2}\footnotesize Degree    & B.Tech EE + M.Tech spec.\ Signal Processing \& ML\\
    \color{ink2}\footnotesize Baselines & 3D Gaussian Splatting (2023) \textperiodcentered\ Mip-Splatting (2024)\\
  \end{tabular}
  \vfill
\end{center}
\clearpage

% ----------------------------------------------------------------- abstract
\chapter*{Abstract}
\addcontentsline{toc}{chapter}{Abstract}
3D Gaussian Splatting builds a radiance field out of several hundred thousand
small translucent ellipsoids, fitted by gradient descent to a set of
photographs. Nothing in that procedure sets a lower bound on how small a
primitive may become. Mip-Splatting supplies one: a scalar variance floor,
$f_k^2 = 0.2\,(d_{\min}/f_{\max})^2$, added in quadrature to all three axes of
every primitive. The argument of this thesis is that the correct bound is
neither scalar nor isotropic, and that both of those properties matter.

The reasoning starts by treating a primitive's position not as something the
optimiser chooses but as something the photographs measure. Each camera that
sees the primitive contributes one noisy observation of where it is, and the
precision of that observation is strongly direction-dependent: a camera pins a
point down well across its line of sight and poorly along it. Accumulating over
all the cameras that saw it gives a $3\times3$ matrix,
\[
  \Lambda_k \;=\; \sum_n w_{nk}\,(J_nW_n)^{\!\top}\,\Sigma_{\mathrm{pix}}^{-1}\,(J_nW_n),
\]
which is the Fisher information of that measurement model at unit weights. Its
eigenvectors are the directions the capture constrained well and badly, and its
eigenvalues say by how much. Cram\'er--Rao motivates reading it as a precision
bound, though the claim actually tested here is weaker: that structure finer
than $1/\sqrt{\lambda_i}$ along eigendirection $i$ is not identifiable from
these particular images.

Three results follow from that expression. The first is that Mip-Splatting's
filter is not a different idea but a special case of it, obtained by keeping one
view, discarding the depth direction and averaging what remains. Its constant
$\sqrt{0.2}$, which the paper introduces without derivation, turns out to be
exactly the product $s\,\sigma_{\mathrm{pix}}$: a confidence radius of $0.447$
pixels. The second is that the anisotropy has a closed form in three regimes.
Within a single view it is already present, falling off as $\cos^2\varphi$ away
from the optical axis. Across two views separated by an angle $\theta$ the
spectrum is triaxial, $\{2,\,2\cos^2(\theta/2),\,2\sin^2(\theta/2)\}$. And
across a \emph{distribution} of views spread over some angular window, as a real capture protocol is, the correct expression is a spherical-cap form
that the two-view result understates by a factor of $\sqrt{2}$. Measured against
that form at each protocol's own width, the observed conditioning of the
$\capAgreeN$ parallax protocols agrees to
$\capAgreeMin$--$\capAgreeMax\,\%$, where the two-view expression falls short by
as much as $\twoViewLowMax\,\%$. The third result is a guarantee, not a formula: if the anisotropic floor is itself floored at Mip-Splatting's own
value, then wherever that value dominates the two methods coincide exactly. On a
well-conditioned benchmark, parity is therefore structural, and an improvement there would indicate a bug.

A second, smaller proposal applies the same reasoning on the sphere of
directions instead of in space, masking the spherical-harmonic coefficients a
capture cannot determine. Those coefficients are $76\,\%$ of the model's
parameters, so the potential saving is large. It does not survive measurement,
and \S\ref{sec:btwo} reports why.

Concurrent work has narrowed what can be claimed for the main idea, and the narrowing is stated here plainly. AAA-Gaussians (2025) already
makes the 3D filter a full $3\times3$, though derived from the view being rendered instead of from the capture. PUP 3D-GS (2025) already builds a
per-primitive second-order spatial sensitivity, though it collapses that to a
scalar and uses it to prune. What remains new is therefore neither the
anisotropy nor the matrix, but where the anisotropy comes from. It comes from the conditioning of the
training capture, fixed once per primitive, and it is kept as a whole spectrum
instead of a summary of one.

\medskip
\keybox{status of this version}{%
Both baselines have been reproduced end to end on free cloud compute.
Mip-Splatting lands within $\stmtArmAMaxDelta$\,dB of its published figures at
all four test scales, over the $\stmtArmAScenes$ scenes the published mean
covers, which meets the numerical-reproduction standard this project set for
itself in advance. So does the gap between the two arms, within
$\stmtGapMaxDelta$\,dB, and that gap is the quantity the whole two-arm
construction exists to produce. The 3DGS column claims nothing, because one
scene will not train on the vanilla arm and the column is therefore a
seven-scene mean standing against an eight-scene published figure.

Two falsification gates were run before any method code existed, since either
could have ended the project: does the resolution floor already bind everywhere,
and is real capture geometry anisotropic enough to matter? Both passed on
measurement.

Several results ran against expectation, and they are reported where they happened. Gate G2 failed, and the first
diagnosis of that failure was itself wrong: the arm labelled 3DGS had kept
Mip-Splatting's screen-space opacity compensation, so it was anti-aliasing better than 3DGS at the very scales the gate was scored
on. Rebuilt and
re-measured, the gate passes on all three criteria. Re-measuring then exposed a second problem, this time in the record itself, in which rows that
a re-run had replaced were never retired and every cell of one column became an
average of two different builds. A third was a denominator: both measured
columns of Table~\ref{tab:stmt} were means over the seven scenes both arms have,
compared against published figures that are means over eight, which inflated
Mip-Splatting's apparent reproduction error fivefold. A fourth was an
explanation. This document had attributed the residual offset to a change in the
densification rule, and said so for months without testing it; measured against
the commit that predates the change, the effect is $\gofMaxEffect$\,dB and is
not consistently signed.

Both proposed methods are implemented and measured. B1 reduces to its baseline
exactly where the theory requires, to $1.4\times10^{-21}$ in the unit test and
to $\methodParityMax$\,dB over $\methodParityScenes$ scenes in full paired
training runs, comfortably below the run-to-run spread. Measurement also amended
the second method before refuting it: an absolute threshold on $\lambda_{\min}$
turns out to have no value that responds gradually across capture regimes, while
the scale-free condition number does.

Whether any of this shows up in a rendered image is Table~\ref{tab:stress}, and
it now has an answer with the shape the derivation demands. On the
$\stressFlatCount$ protocols where the resolution floor still dominates, the
estimation term exceeds it on zero primitives, the filter comes out isotropic to
two decimal places, and B1 differs from Mip-Splatting by between
$\stressFlatMin$ and $\stressFlatMax$\,dB. On the sixth, a three-camera pencil
spanning $29^\circ$, the estimation term exceeds the floor on $93\,\%$ of
primitives, the filter becomes anisotropic at $1.47\times$, and B1 leads by
$\stressBindDelta$\,dB, on the mean and on every one of the eight paired cells.
The step falls exactly where the estimation term overtakes the floor, as predicted and recorded
beforehand.

Two things that result is not. The pencil reconstructs at $14$\,dB against the
full orbit's $34$, so what has been established is that the mechanism is the one
claimed, not that the method is useful at that operating point. And eight
consistent pairings are still eight cells of a single run on two scenes.

\medskip
Two results are negative, and they belong here.
\textbf{B2 is refuted.} Compared against a control that has B2's own
nested-degree structure and B2's own budget, so that only the criterion differs,
magnitude pruning wins by up to $\btwoBlockedGapMax$\,dB. Chapter~\ref{ch:method2}
had set the standard that a criterion unable to beat magnitude pruning at
matched size is not worth the machinery, and by that standard it fails. An
earlier and much larger-looking margin, measured against a control with more
freedom than B2 has, turned out to be mostly a verdict on the shape of the mask, leaving the
criterion itself barely implicated; that correction is reported too. \textbf{B1 costs
$2.8\times$ the peak memory and $1.32\times$ the training time} of the baseline
it provably equals on a well-conditioned capture. On such a capture it should
not be used, and this thesis does not recommend it there.}
\clearpage

\chapter*{Current Status}
\addcontentsline{toc}{chapter}{Current Status}

\paragraph{Reproduction.} Both baselines are reproduced end to end on free cloud
compute. Mip-Splatting lands within $\stmtArmAMaxDelta$\,dB of its published
figures at all four test scales, over the $\stmtArmAScenes$ scenes the published
mean covers, which meets the numerical-reproduction standard set in advance. So
does the gap between the two arms, within $\stmtGapMaxDelta$\,dB, and that gap
is the quantity the two-arm construction exists to produce. The 3DGS column
claims nothing: one scene will not train on that arm, leaving a seven-scene mean
standing against an eight-scene published figure.

\paragraph{Falsification Gates.} Two were run before any method code existed,
since either could have ended the project: does the resolution floor already
bind everywhere, and is real capture geometry anisotropic enough to matter? Both
passed on measurement.

\paragraph{The Anisotropic Filter.} Implemented, and it reduces to its baseline
where the theory requires: to $1.4\times10^{-21}$ in the unit test, and to
$\methodParityMax$\,dB over $\methodParityScenes$ scenes in paired training
runs, below the run-to-run spread. It costs $2.8\times$ the peak memory and
$1.32\times$ the training time of the baseline it equals on a well-conditioned
capture, so on such a capture it should not be used.

\paragraph{The Stress Suite.} On the $\stressFlatCount$ protocols where the
resolution floor dominates, the estimation term exceeds it on zero primitives,
the filter is isotropic to two decimals, and the difference from Mip-Splatting
runs between $\stressFlatMin$ and $\stressFlatMax$\,dB. On the sixth, a
three-camera pencil spanning $29^\circ$, the estimation term exceeds the floor
on $93\,\%$ of primitives, the filter becomes anisotropic at $1.47\times$, and
the anisotropic filter leads by $\stressBindDelta$\,dB on the mean and on every
one of eight paired cells. The step falls where the estimation term overtakes
the floor, as predicted and recorded beforehand. Two qualifications: the pencil
reconstructs at $14$\,dB against the full orbit's $34$, so what is established is
the mechanism and not the method's usefulness at that operating point; and eight
consistent pairings are eight cells of one run on two scenes.

\paragraph{The Angular Mask.} Refuted. Against a control carrying its own
nested-degree structure and its own budget, so that only the criterion differs,
magnitude pruning wins by up to $\btwoBlockedGapMax$\,dB.
Chapter~\ref{ch:method2} set the standard that a criterion unable to beat
magnitude pruning at matched size is not worth the machinery, and by that
standard it fails.

\paragraph{Corrections Made During The Work.} Four, each reported where it
happened. The arm labelled 3DGS had kept Mip-Splatting's screen-space opacity
compensation, so it was anti-aliasing better than 3DGS at the scales the gate
was scored on; rebuilt and re-measured, the gate passes on all three criteria.
Rows that a re-run had replaced were never retired, so every cell of one column
became an average of two builds. Both measured columns of
Table~\ref{tab:stmt} were means over seven scenes compared against published
figures that are means over eight, which inflated the apparent reproduction
error fivefold. And the residual offset had been attributed to a change in the
densification rule without being tested; measured against the commit that
predates that change, the effect is $\gofMaxEffect$\,dB and is not consistently
signed.
\clearpage


\tableofcontents
\clearpage

\part{Problem}
\chapter{Capture-Limited Resolution}
\label{ch:intro}

Suppose you are handed eighty photographs of an object and asked to produce a
picture of it from a viewpoint nobody stood at. This is the problem a radiance field solves. Before any of the machinery, consider what the photographs can tell you, and
where they run out.

They sample the scene, and they sample it unevenly. A surface close to the
cameras is covered by many more pixels than one far away. Some parts of the
object are seen from many directions, others only from a narrow range, and a few
may be seen from a single grazing angle or not at all. Different cameras may not
even have the same focal length, so their pixels do not correspond to the same
patch of surface. Whatever you reconstruct from such a set, there is a limit on
how fine the detail can be, and that limit belongs to the capture. It follows the photographs into whatever
representation stores the answer, and no choice of representation lifts it.

NeRF acts on this directly. Its positional encoding uses $L$ frequency bands,
and $L$ is chosen from, in the authors' words, ``the maximum frequency present
in the sampled input images''. The quantity is the right one; the instrument is
one number, picked by hand, applied to an entire scene.

3D Gaussian Splatting drops even that. Its authors state that no regularisation
is applied, and a primitive may therefore shrink without bound: nothing in the
objective opposes it. As long as the
model is evaluated at the resolution it was fitted at, this does no visible harm, which is why it went unnoticed for so long. Change the
resolution and the damage appears at once. Section~\ref{sec:cost} puts a number
on it: $15.64$\,dB in the published figures, and \StmtArmBDrop\,dB in this
project's own reproduction.

Mip-Splatting restored a limit, and the restoration works; the numbers in
Chapter~\ref{ch:results} confirm it comfortably. But the limit it restores is a
single scalar per primitive, added in quadrature to all three axes of the
covariance. This thesis asks what that scalar ought to have been, and answers that the right
object is a matrix.

\section{The Claim}

\keybox{thesis}{%
``The maximum frequency present in the sampled input images'' varies from place
to place and from direction to direction. It falls with depth; it differs across
the ray from along it; and it varies over the sphere of viewing directions. The accumulated sampling Fisher information is its
local, directional generalisation. Floored at the per-view Nyquist limit, it
reduces to Mip-Splatting exactly where the capture is well conditioned, and
departs from it precisely where the capture is ill conditioned.}

\section{The Signal-Processing Problem}

Stripped of its graphics vocabulary, the problem is reconstruction from
non-uniform samples without a matched prefilter --- a classical
signal-processing failure mode. The objects involved are familiar ones.

\begin{table}[htbp]
  \centering
  \caption{The same objects, in two vocabularies.}
  \small
  \begin{tabularx}{\linewidth}{L{0.91}L{1.09}}
    \toprule
    graphics term & signal-processing object \\
    \midrule
    Anti-aliasing filter (2D Mip filter) & Reconstruction prefilter matched to the display sampling rate \\
    3D smoothing filter & Prefilter matched to the \emph{acquisition} sampling rate \\
    Primitive covariance $\Sigma_k$ & Support of a basis function in a non-orthogonal expansion \\
    ``Too many small Gaussians'' & Over-parameterisation relative to the information in the measurements \\
    Floaters and depth artefacts & Fitting components along directions the data does not constrain \\
    Spherical-harmonic view dependence & Basis expansion whose Gram matrix may be rank-deficient \\
    \bottomrule
  \end{tabularx}
\end{table}

The contribution is to compute the acquisition-side prefilter from the conditioning of the capture instead of choosing it by hand, and to keep it anisotropic
because the information is. AAA-Gaussians (2025) already has an anisotropic 3D filter, derived from the view
being rendered, and PUP 3D-GS (2025) already builds a per-primitive second-order
information matrix in order to prune. Section~\ref{sec:relatedwork} states the line of
separation precisely and Chapter~\ref{ch:discussion} returns to it; the short
form is that what is new here is the \emph{source} of the anisotropy and the
retention of its \emph{spectrum}.

\section{Contributions}

\begin{enumerate}[leftmargin=1.4em,itemsep=3pt]
  \item \textbf{A band-limit derived from the capture.} The finest structure a
        primitive may carry, per direction, computed from an observability
        matrix accumulated over the cameras that saw it, using the same Jacobian
        the renderer already forms
        (\S\ref{sec:jacobian}--\S\ref{sec:accumulation}). The limit is a
        property of the capture, fixed once per primitive, and its whole
        spectrum is kept.

  \item \textbf{Mip-Splatting's filter identified as a special case.} It is the
        single-view, lateral-only, isotropised reduction of that expression, and
        its undocumented constant $\sqrt{0.2}$ is $s\,\sigma_{\mathrm{pix}}$: a
        confidence radius of $0.447$ pixels (\S\ref{sec:proposition1}).

  \item \textbf{Closed forms for the anisotropy in three regimes.} Within one
        view, $\lambda_2/\lambda_1 = \cos^2\varphi$ off-axis
        (\S\ref{sec:offaxis}). Across two views, the triaxial spectrum
        $\{2,\,2\cos^2(\theta/2),\,2\sin^2(\theta/2)\}$
        (\S\ref{sec:anisotropy}). Across a distribution of views, the
        spherical-cap form (\S\ref{sec:capprediction}), which a real protocol
        obeys and which the two-view expression understates by $\sqrt{2}$.
        Together they predict where the method acts and where it cannot.

  \item \textbf{An exact-reduction theorem.} Flooring the anisotropic limit at
        Mip-Splatting's own value makes the two coincide wherever that value
        dominates, so parity on a well-conditioned benchmark is structural
        (\S\ref{sec:proposition2}).

  \item \textbf{A criterion for angular identifiability of spherical
        harmonics}, which are $76\,\%$ of the model, and a controlled test of it
        against magnitude pruning at matched structure and matched budget
        (Chapter~\ref{ch:method2}, \S\ref{sec:btwo}).

  \item \textbf{An executed reproduction of both baselines} on free compute,
        with twelve pre-flight assertions and the per-scale recovery the
        published evaluation does not provide (Chapters~\ref{ch:design}
        and~\ref{ch:results}).

  \item \textbf{A measured cost for the method}, paired in one session on one
        card: $2.8\times$ peak memory and $1.32\times$ training time for a model
        whose primitive count matches the baseline's to $0.2\,\%$
        (\S\ref{sec:methodcost}).

  \item \textbf{Checks that catch measurement errors which read as plausible.}
        A build that refuses any table whose cells span two code versions, a
        test that holds the thesis, the figures and the tabulated results to the
        same numbers, and a record that keeps every superseded run. Four errors
        in this project's own measurement were found this way, each reported
        with the correction it forced (Appendix~\ref{app:defects}).
\end{enumerate}

\input{ch02-splatting-and-band-limits}
\input{ch03-scalar-limitations}

\part{Method}
\chapter{An anisotropic band-limit}
\label{ch:method1}

\keybox{the idea}{%
Think of each blob's position as something the photographs \emph{measure}. Every training photograph that sees the blob is a
measurement of where it is, and every measurement has noise --- roughly a pixel's
worth, in the image plane.

Now ask the standard estimation question: given all those measurements, how precisely do you know the position? The answer varies with direction. A camera
localises a point very well within its image plane and poorly along its own
viewing direction; a second camera at a different angle fixes the first camera's
weak direction. Accumulate over all the cameras that saw the blob and you get a
$3\times3$ matrix --- the familiar Fisher information of the measurement model, written $\Lambda_k$
throughout. Its eigenvectors are the directions the capture
constrains well and badly; its eigenvalues say how well.

The proposal is then almost forced: \emph{a primitive should not be finer, in any
direction, than the capture can determine it in that direction}. Invert the
matrix, and you have a covariance floor with the right shape. Then floor that at the per-view Nyquist limit, since no amount of accumulated
information lets you resolve structure finer than a pixel, and you have the
method.

Three things follow, and this chapter proves them. Mip-Splatting's scalar filter
is what you get by collapsing this matrix to a single view and throwing
away direction. Its mysterious constant $\sqrt{0.2}$ turns out to be a confidence
radius of $0.447$ pixels. And when the Nyquist floor dominates, which on a normal benchmark capture it
always does, the new method becomes the old one exactly. That last point is why the standard benchmarks in
Chapter~\ref{ch:results} show no improvement, and why a gain there would have indicated a bug.}

\begin{figure}[htbp]
  \centering
  \includegraphics[width=0.97\linewidth]{figures/d3-one-camera-two-cameras.pdf}
  \caption{Why the limit has to be a matrix. A single camera fixes a point well
  across its line of sight and badly along it, so the region consistent with its
  measurement is a long thin cigar pointing back at the lens --- the familiar
  fact that one photograph gives you direction but not distance. A second camera from elsewhere is precise where the first is vague, and the region
  consistent with both is compact. Accumulating this over every camera that saw
  the primitive is what $\Lambda_k$ does, and its eigenvectors are these axes.
  A scalar limit cannot distinguish the long axis from the short one; that is
  the whole of the argument.}
  \label{fig:cameras}
\end{figure}

\section{The estimation problem}
Regard the position $p_k$ of a primitive as a parameter to be estimated from
image measurements. Camera $n$ observes the primitive's projection at
\begin{equation}
  u_n \;=\; \pi_n(p_k) + \varepsilon_n,\qquad
  \varepsilon_n \sim \mathcal{N}(0,\Sigma_{\mathrm{pix}}),\quad
  \Sigma_{\mathrm{pix}} = \sigma_{\mathrm{pix}}^2 I_2 ,
  \label{eq:measurement}
\end{equation}
where $\sigma_{\mathrm{pix}}$ is the localisation noise of an image feature, in
pixels. The log-likelihood of the measurement from view $n$ is, up to constants,
$-\tfrac12 (u_n-\pi_n(p))^{\!\top}\Sigma_{\mathrm{pix}}^{-1}(u_n-\pi_n(p))$.

\section{The measurement Jacobian}
\label{sec:jacobian}
The Fisher information carried by view $n$ about $p$ is
\begin{equation}
  I_n \;=\; \Big(\frac{\partial \pi_n}{\partial p}\Big)^{\!\top}
            \Sigma_{\mathrm{pix}}^{-1}
            \Big(\frac{\partial \pi_n}{\partial p}\Big),
  \qquad \frac{\partial \pi_n}{\partial p} = J_n W_n .
  \label{eq:fisher-view}
\end{equation}
The derivative of image position with respect to world position is the EWA pair
of Equation~\eqref{eq:ewa}. \emph{No new geometry is introduced}: the quantity
the renderer forms to project a covariance is the same quantity estimation
theory requires to bound a position. This identification is the technical core
of the thesis, and it is what makes the method cheap.

\paragraph{Rank.} $J_nW_n \in \mathbb{R}^{2\times3}$, so $I_n$ has rank at most
$2$. A single view constrains the two directions transverse to the ray and says
nothing about depth. For a pinhole camera the eigenvalues of $I_n$ are exactly
$(0,\ (f/z)^2/\sigma_{\mathrm{pix}}^2,\ (f/z)^2/\sigma_{\mathrm{pix}}^2)$
(Appendix~\ref{app:rank}). Depth information exists only through parallax.

\section{The observability matrix}
\label{sec:accumulation}
Summing over views, weighted by each view's visibility and blending contribution $w_{nk}$, which the rasteriser already accumulates,
\begin{equation}
  \Lambda_k \;=\; \sum_n w_{nk}\,(J_nW_n)^{\!\top}\Sigma_{\mathrm{pix}}^{-1}(J_nW_n)
  \;=\; V\,\mathrm{diag}(\lambda_1,\lambda_2,\lambda_3)\,V^{\!\top}.
  \label{eq:lambda}
\end{equation}

\paragraph{The weight, defined.} $w_{nk} = T_{nk}\,\alpha_k \in [0,1]$ is the
primitive's blending contribution in view $n$: the transmittance accumulated in
front of it times its own opacity, which is the coefficient the rasteriser already multiplies its colour by. It is zero for a primitive the view
does not see and near one for an unoccluded opaque one. This weighting is a \emph{heuristic} down-weighting of occluded views, standing outside any
likelihood:
a view that cannot see a primitive carries no information about it, and $w_{nk}$
is the cheapest available statement of that fact. It is stated as a heuristic
because calling it anything more would be a claim the derivation does not
support.

\paragraph{What the matrix is, precisely.} With unit weights,
Equation~\eqref{eq:lambda} \emph{is} the Fisher information of the measurement
model in Equation~\eqref{eq:measurement}. With the weights it departs from that, so the honest name for $\Lambda_k$ is an
\textbf{observability matrix} --- a
view-weighted accumulation of the per-view measurement information, coinciding
with the Fisher information exactly when $w_{nk} \equiv 1$. Under the
point-feature approximation of Equation~\eqref{eq:measurement},
Equation~\eqref{eq:lambda} is also the Gauss--Newton approximation to the Hessian
of the photometric cost in $p_k$: the residual is linearised through $J_nW_n$ and second-derivative terms are dropped, as every bundle adjustment does and for
the same reason.

\paragraph{Why this matters for the claim.} The Cram\'er--Rao bound
$\mathrm{Cov}(\hat p_k) \succeq \Lambda_k^{-1}$ holds for an unbiased estimator
of the measurement model as written. A splatting optimiser is not that estimator:
it fits a photometric loss over a densifying set of correlated primitives, not
independent point features with Gaussian localisation noise. Cram\'er--Rao is
therefore kept as the \emph{motivation} for the form of Equation~\eqref{eq:filtervar}: it says why
the inverse of an information matrix is the right shape for a resolution limit.
The claim the experiments actually test is weaker, and sufficient: along eigendirection $i$ the capture constrains position
to order $1/\sqrt{\lambda_i}$, so structure finer than that is not identifiable
from these images.

\keybox{principle}{%
A primitive should not be narrower, along a given direction, than the
uncertainty with which the capture determines its position along that direction.
Structure finer than $s/\sqrt{\lambda_i}$ is not identifiable from the
measurements; fitting it is fitting noise.}

\section{The covariance floor}
Hence a per-direction variance, floored at Mip-Splatting's own value and capped
at a scene-scale fraction:
\begin{equation}
  \sigma_i^2 \;=\; \min\!\Big(\max\big(f_k^2,\ s^2/\lambda_i\big),\ (\beta z_{\min,k})^2\Big),
  \label{eq:filtervar}
\end{equation}
\begin{equation}
  \Sigma_{\mathrm{filt}} = V\,\mathrm{diag}(\sigma_1^2,\sigma_2^2,\sigma_3^2)\,V^{\!\top};
  \quad \Sigma_k \leftarrow \Sigma_k + \Sigma_{\mathrm{filt}};
  \quad \alpha_k \leftarrow \alpha_k\sqrt{|\Sigma_k|/|\Sigma_k+\Sigma_{\mathrm{filt}}|}.
  \label{eq:filter}
\end{equation}
Two hyperparameters are introduced and none removed. $\beta$ caps the filter in
degenerate directions where $\lambda_i \to 0$; without it an unconstrained
direction would produce an unbounded primitive. The opacity compensation is
Mip-Splatting's own, unchanged, and Appendix~\ref{app:mass} shows it carries over
to an anisotropic $\Sigma_{\mathrm{filt}}$ without modification.

\paragraph{$s$ and $\sigma_{\mathrm{pix}}$, disentangled.} These two appear only
as the product $s\,\sigma_{\mathrm{pix}}$, so quoting a default for one and a
value for the product over-determines them. The convention used throughout, and
the one the implementation follows:
\begin{itemize}[leftmargin=1.4em,itemsep=2pt]
  \item $\sigma_{\mathrm{pix}} = \sqrt{0.2} = 0.447$\,px is \emph{inherited} from Mip-Splatting's own constant. Proposition~\ref{prop:constant} shows it is the localisation
        noise implicit in Mip-Splatting's constant, and taking it verbatim is
        what makes Proposition~\ref{prop:reduction} exact instead of merely close.
  \item $s$ is the dimensionless confidence multiplier, default $\mathbf{1}$
        (one standard deviation). It is the only free knob of the pair, and
        every result in this thesis is at $s=1$.
\end{itemize}
So $s\,\sigma_{\mathrm{pix}} = 0.447$\,px at the default, and a reader changing
$s$ changes the confidence level and nothing else.

\paragraph{Consistency of the floor and the cap.} The floor must never exceed the
cap, or Equation~\eqref{eq:filtervar} would be ill-posed. Requiring
$f_k \le \beta z_{\min,k}$ and substituting Equation~\eqref{eq:mipfilter} gives
$\beta \ge \sqrt{0.2}/f_{\max}$. For $f_{\max}\approx1000$\,px this is
$\beta \ge 4.5\times10^{-4}$; the default $\beta=0.01$ carries $22\times$
headroom.

\section{Mip-Splatting as a special case}
\label{sec:proposition1}
Take a single camera with focal length $f$ viewing a primitive at depth $z$ on
the optical axis. From \S\ref{sec:jacobian} the lateral eigenvalue is
$\lambda_{\mathrm{lat}} = (f/z)^2/\sigma_{\mathrm{pix}}^2$, so the
estimation-limited lateral scale is $s/\sqrt{\lambda_{\mathrm{lat}}} =
s\,\sigma_{\mathrm{pix}}\,z/f$.

\begin{proposition}[Identification of the constant]
\label{prop:constant}
Mip-Splatting's 3D smoothing filter is exactly Equation~\eqref{eq:filtervar}
evaluated for a single view, restricted to the lateral directions, and applied
isotropically --- provided
\[ s\cdot\sigma_{\mathrm{pix}} \;=\; \sqrt{0.2} \;\approx\; 0.447\ \text{pixels}. \]
\end{proposition}

Their constant therefore carries a meaning: a sub-pixel localisation confidence
of $0.447$\,px, slightly under half a pixel. The present method keeps
that calibration and generalises the geometry around it. This is the strongest
single argument in the thesis: it converts the relationship between the proposed
method and its baseline from ``similar in spirit'' into ``the baseline is a
corollary'', and it explains a magic number in a published paper.
Appendix~\ref{app:prop1} gives the derivation in full.

\section{Off-axis anisotropy}
\label{sec:offaxis}

Before parallax enters at all, there is anisotropy the scalar filter discards.
Section~\ref{sec:jacobian} quotes the single-view eigenvalues
$(0, (f/z)^2, (f/z)^2)/\sigma_{\mathrm{pix}}^2$, and that is exact --- \emph{on
the optical axis}. The convenient rank-1 form
$(f/z)^2(I - dd^{\!\top})/\sigma_{\mathrm{pix}}^2$ suggests otherwise, because
$I - dd^{\!\top}$ is a projector and a projector's non-zero eigenvalues are equal
by construction, whatever $d$ is. Off-axis, the true $J_n^{\!\top}J_n$ departs from a projector, and the two
lateral eigenvalues separate.

Writing $a = (x/z,\,y/z)$ for the normalised off-axis displacement,
$J = (f/z)B$ with $B = \big[\begin{smallmatrix}1&0&-a_1\\0&1&-a_2\end{smallmatrix}\big]$,
so $BB^{\!\top}$ has trace $2+|a|^2$ and determinant $1+|a|^2$. Its eigenvalues
are therefore $1+|a|^2$ and $1$, giving
\begin{equation}
  \frac{\lambda_2}{\lambda_1} \;=\; \frac{1}{1+|a|^2} \;=\; \cos^2\varphi,
  \qquad \varphi = \arctan|a| \ \text{the off-axis angle.}
  \label{eq:offaxis}
\end{equation}

\input{generated/table-geometry}

At the corner of a $50^\circ$-field camera $\varphi$ approaches $25^\circ$ and
the two lateral eigenvalues already differ by $20\,\%$; the sampled poses of
Table~\ref{tab:offaxis} run from $0.93$ down to $0.47$. So \textbf{perspective
contributes anisotropy independently of parallax}, and a scalar filter discards
it before the parallax argument of \S\ref{sec:anisotropy} has even begun. The
implementation uses the full $J_nW_n$ and never the rank-1 shortcut, so this term survives into the filter, where the shortcut would discard it for
good; the shortcut appears in this document
only where it makes an argument shorter, and Equation~\eqref{eq:offaxis} records
what it costs. This is also the most likely explanation for the residual in
\S\ref{sec:instrumentresults}: the grazing protocol, which is a full orbit tilted
towards obliquity, measures more anisotropy than any parallax-only model
predicts.

\section{The two-view spectrum}
\label{sec:anisotropy}
Take the simplest case: two cameras at the same distance $z$, placed symmetrically about the primitive, so that the angle
between their lines of sight is $\theta$. Evaluating
Equation~\eqref{eq:lambda} for this arrangement (the algebra is in
Appendix~\ref{app:twoview}) gives $\Lambda \propto 2I - d_1d_1^{\!\top} -
d_2d_2^{\!\top}$, and this has the exact spectrum
$\{2,\ 2\cos^2(\theta/2),\ 2\sin^2(\theta/2)\}$.

These are three distinct numbers. The right filter is therefore triaxial, with a different limit along each of
three axes: the lateral
direction lying in the plane of the two rays is degraded as well, and only the
direction perpendicular to that plane is fully constrained. Taking the extreme
pair of eigenvalues,
\begin{equation}
  \frac{\lambda_3}{\lambda_1} = \sin^2(\theta/2),
  \qquad
  \frac{\sigma_{\mathrm{depth}}}{\sigma_{\mathrm{lat}}} = \frac{1}{\sin(\theta/2)}
  \qquad (\theta \le 90^\circ).
  \label{eq:anisotropy}
\end{equation}
This is exact at every angle --- verified against pinhole Jacobians to
$4\times10^{-16}$ over $\theta = 2^\circ$ to $179^\circ$
(Table~\ref{tab:twoviewspectrum}). The restriction $\theta \le 90^\circ$ earns its place: past $90^\circ$ the depth direction becomes better constrained than the
in-plane lateral one and the two exchange roles, so the ratio of smallest to
largest is $\min(\sin^2(\theta/2), \cos^2(\theta/2))$. At $\theta = 120^\circ$
the naive form reads $0.75$ where the truth is $0.25$ --- a factor of three, in
the optimistic direction. Every protocol this thesis stresses lies well below
$90^\circ$, and the full orbit lies so far above it that the ratio is $\approx 1$
either way, so the restriction costs nothing here and guards against a real error elsewhere.

\subsection{The spherical-cap correction}
\label{sec:capprediction}

Equation~\eqref{eq:anisotropy} is exact for \emph{two} cameras. A capture
protocol is ten or a hundred, spread over an angular window, and $\Lambda$ is the
sum over that spread --- a different number. Reading the two-view formula off a protocol's measured span is defensible only
for a protocol that really is two cameras, and every protocol here has more.

For cameras distributed uniformly over a spherical cap of half-angle
$\alpha = \theta/2$ about the mean viewing direction, $\cos t$ is uniform on
$[\cos\alpha, 1]$, so with $c = \cos\alpha$ and $E = \mathbb{E}[d_z^2] =
(1+c+c^2)/3$,
\begin{equation}
  \frac{\lambda_3}{\lambda_1} \;=\; \frac{2(1-E)}{1+E},
  \qquad
  \frac{\sigma_{\mathrm{depth}}}{\sigma_{\mathrm{lat}}}
  \;=\; \sqrt{\frac{1+E}{2(1-E)}} .
  \label{eq:cap}
\end{equation}

\keybox{the two-view form against a real capture}{%
As $\alpha \to 0$, Equation~\eqref{eq:cap} tends to $\alpha^2/2$ while
$\sin^2\alpha \to \alpha^2$: the cap's $\lambda_3/\lambda_1$ is
\textbf{half} the two-view value, because a cap's cameras are spread across the window instead of sitting at its two ends. A smaller $\lambda_3/\lambda_1$ is a
\emph{worse}-conditioned capture, so Equation~\eqref{eq:anisotropy} applied to a
protocol \emph{understates} its depth anisotropy by $\sqrt{2}$ in $\sigma$. The
direction matters: the conditioning instrument measuring more anisotropy than the two-view form predicts
is the model working, not failing, and reporting it against the two-view number
would have made a confirmation look like a discrepancy. Every prediction quoted
against a measured span in this thesis therefore uses Equation~\eqref{eq:cap}.}

\begin{figure}[htbp]
  \centering
  \includegraphics[width=\linewidth]{figures/fig5-fisher-geometry.pdf}
  \caption{The mechanism. Cross-sections of the two filters in the plane of the
  cameras, normalised to the same lateral extent because the floor makes them
  agree there by construction, so the whole visible difference is along depth.
  The two panels are the spans the \texttt{cone} and \texttt{arc} protocols
  actually realise ($20^\circ$ and $80^\circ$), and the elongation is
  Equation~\eqref{eq:cap} at those spans --- $8.1\times$ and $2.0\times$. On a
  full orbit the two ellipses coincide to within the line width, which is Proposition~\ref{prop:reduction} drawn instead of stated. Geometry computed from real pinhole Jacobians.}
  \label{fig:fisher}
\end{figure}

\begin{figure}[htbp]
  \centering
  \includegraphics[width=0.86\linewidth]{figures/fig6-anisotropy.pdf}
  \caption{Where the method has room, and which prediction to use. Mip-Splatting
  assumes the ratio is $1$ at every span; the gap between that line and the
  curve is the quantity this thesis exploits. Both forms are drawn because the
  difference between them is a correction this project had to make: the dashed
  curve is Equation~\eqref{eq:anisotropy}, exact for two cameras, and the solid
  one is Equation~\eqref{eq:cap} for cameras distributed over the span, as a protocol is. The points are the conditioning instrument measured on eight scenes
  (\S\ref{sec:instrumentresults}), and they sit above the two-view curve at
  every protocol --- plotted against that curve alone, a confirmation would read
  as a discrepancy.}
  \label{fig:anisotropy}
\end{figure}

\begin{table}[htbp]
  \centering
  \caption{What Equation~\eqref{eq:cap} predicts, at the spans the protocols of
  \S\ref{sec:stress} actually realise. A set of falsifiable predictions, made
  before the experiments, and the reason the stress suite exists. A method that helped everywhere would contradict its own derivation.}
  \label{tab:regimes}
  \small
  \begin{tabular}{lccl}
    \toprule
    capture & span & $\sigma_D/\sigma_L$ & expectation \\
    \midrule
    full orbit (standard benchmarks) & $179^\circ$ & $1.00\times$ & parity --- the floor dominates \\
    grazing orbit & $168^\circ$ & $1.04\times$ & parity from parallax alone \\
    one-sided arc & $80^\circ$ & $2.03\times$ & gain along depth \\
    low-parallax cone & $20^\circ$ & $8.10\times$ & largest gain \\
    single view & $0^\circ$ & $\infty$ & cap $\beta$ binds; scene-scale prior \\
    \bottomrule
  \end{tabular}
  \par\vspace{2pt}\footnotesize Spans are the medians the built protocols
  measured (Table~\ref{tab:instruments}), so the prediction and the measurement
  are evaluated at the same geometry. Quoting Equation~\eqref{eq:anisotropy}
  here instead would give $1.00$, $1.01$, $1.56$ and $5.76\times$ --- uniformly
  low, for the reason \S\ref{sec:capprediction} gives.
\end{table}

\section{Exact reduction}
\label{sec:proposition2}

\begin{figure}[htbp]
  \centering
  \includegraphics[width=0.97\linewidth]{figures/d5-floor-vs-estimate.pdf}
  \caption{Why a benchmark result of ``no change'' is the correct outcome. Two limits are in play and the method takes whichever
  is larger. The flat blue line is the resolution floor: no capture, however
  generous, lets you resolve structure finer than a pixel. The rising orange
  curve is the estimation limit of this chapter, which grows as the cameras
  crowd into a narrower window. On a standard benchmark, meaning a hundred cameras around a full circle, the
  floor sits far above the estimation limit, the
  method returns the floor, and it \emph{is} Mip-Splatting, exactly and provably
  (Proposition~\ref{prop:reduction}). Only when the capture narrows past the
  crossing does anything change. The four marks are the protocols of
  \S\ref{sec:stress} at their measured widths; the crossing falls between the
  cone and the pencil, which is where the measurement of \S\ref{sec:table3} puts
  it.}
  \label{fig:floorvsest}
\end{figure}

\begin{proposition}[Exact reduction]
\label{prop:reduction}
If $s^2/\lambda_i \le f_k^2$ for every $i$, then Equation~\eqref{eq:filter} is
identical to Mip-Splatting's Equation~\eqref{eq:mipfilter}.
\end{proposition}

\begin{proof}
Under the hypothesis, and given $\beta \ge \sqrt{0.2}/f_{\max}$ so the cap does
not bind, Equation~\eqref{eq:filtervar} gives $\sigma_i^2 = f_k^2$ for all $i$.
Then $\Sigma_{\mathrm{filt}} = V(f_k^2 I)V^{\!\top} = f_k^2 VV^{\!\top} = f_k^2 I$,
since $V$ is orthogonal. The opacity compensation of Equation~\eqref{eq:filter}
then coincides with Equation~\eqref{eq:opacitycomp}.
\end{proof}

\paragraph{Consequence.} The filter never smooths \emph{less} than the baseline
in any direction. Parity on well-conditioned benchmarks is therefore a
structural property, not an empirical hope, and a failure to achieve it is a
bug, testable to numerical precision.

This proposition is implemented as a unit test: with the accumulation disabled the model must reproduce Mip-Splatting to $10^{-6}$,
and with the filter zeroed entirely it must reproduce 3DGS. The second of these
is already in force. It \emph{is} the 3DGS-arm construction of \S\ref{sec:onecodebase}, and Chapter~\ref{ch:results} reports it verified on the saved models.

\section{Implementation cost}
\label{sec:implcost}

Two consequences follow from $\Sigma_{\mathrm{filt}}$ being a matrix. Both belong to the implementation.

\paragraph{The full $3\times3$ must reach the rasteriser.} Mip-Splatting can add
its filter by updating each scale as $s_i \leftarrow \sqrt{s_i^2 + f_k^2}$,
because $f_k^2 I$ commutes with $\Sigma_k$ and shares its eigenvectors. An
anisotropic $\Sigma_{\mathrm{filt}}$ does not commute with $\Sigma_k$ in general:
$V$ diagonalises $\Lambda_k$ and $R$ diagonalises $\Sigma_k$, and nothing aligns
them. So $\Sigma_k + \Sigma_{\mathrm{filt}}$ is not expressible as a per-axis
update of the scale vector, and the summed covariance has to be formed in full
and passed through the rasteriser's precomputed-covariance input. The implementation does that. Where the filter comes out isotropic, which Proposition~\ref{prop:reduction}
makes the common case, it takes Mip-Splatting's own scalar path instead, so parity holds in the implementation as well as in the algebra: the two take
the same code path.

\paragraph{The eigendecomposition stays off the backward path.} $\Lambda_k$ is
computed under \texttt{no\_grad} and held constant between recomputations,
exactly as \texttt{filter\_3D} is. The gradient therefore never passes through
it, and the question of differentiating a possibly degenerate eigenvector basis
never arises, though it is where a naive implementation of
this idea would have failed. The filter is a piecewise-constant
band-limit that the optimiser sees as a fixed additive covariance, held constant throughout
training.

\section{Arithmetic cost}
$\Lambda_k$ is a sum of $N$ outer products of a $2\times3$ matrix, followed by a
$3\times3$ symmetric eigendecomposition per primitive. It is recomputed on the
same schedule Mip-Splatting already uses for \texttt{filter\_3D} --- every 100
iterations, roughly 290 times across a 30\,000-iteration run, and never at
render time. Rendering speed is therefore unchanged by construction, which is
the answer to the first question a graphics audience will ask, and \S\ref{sec:fps} measures it instead of leaving it as an argument.

\chapter{Angular identifiability of spherical harmonics}
\label{ch:method2}

\begin{figure}[htbp]
  \centering
  \includegraphics[width=0.99\linewidth]{figures/d7-sh-coverage.pdf}
  \caption{The argument of this chapter, before the algebra. A primitive's
  colour is allowed to depend on the direction it is viewed from, and that
  dependence is stored as a short series in the viewing direction. A
  degree-$\ell$ term oscillates $\ell$ times around the sphere, so telling the
  degrees apart requires seeing the primitive from enough different directions.
  Seen from all around (left) that is easy. Seen from an eleven-degree window
  (right) it is impossible: rescaled to the window, all three curves are the
  same gentle arc, and no fit can say which of them the data was asking for.
  Fitting all three anyway does not produce an error; it produces three
  coefficients that trade off against one another arbitrarily.}
  \label{fig:shcoverage}
\end{figure}

The same argument applies on the sphere of directions. Each primitive carries
view-dependent colour as spherical harmonics to degree $3$: $16$ coefficients per
channel, $48$ floats, against $11$ for position, scale, rotation and opacity
combined. Non-DC spherical harmonics are $45$ of $59$ floats --- $76\,\%$ of the
model. They are fitted from however many directions happen to see the primitive,
which may be very few.

Let $y(d)$ be the vector of SH basis functions evaluated at direction $d$. The
angular Gram matrix accumulated over views is
\begin{equation}
  A_k \;=\; \sum_n w_{nk}\, y(d_{nk})\, y(d_{nk})^{\!\top}.
  \label{eq:gram}
\end{equation}
A coefficient block is determinable only where the corresponding block of $A_k$
is well conditioned. The criterion is to retain degree $\ell$ for primitive $k$
only where
\begin{equation}
  \ell_{\max}(k) \;=\; \max\big\{\, \ell \;:\; \lambda_{\min}\big(A_k(\ell)\big) > \tau \,\big\},
  \label{eq:lmax}
\end{equation}
and to mask higher degrees to zero. A primitive seen from three directions
cannot support degree $3$; fitting it produces coefficients that encode noise
and, worse, produce visible artefacts under novel views that fall outside the
observed cone.

\keybox{why this is a separate contribution}{%
The angular mask is independent of the anisotropic filter. It survives even if the floor-occupancy diagnostic of
\S\ref{sec:instruments} kills the positional band-limit: the two share a principle --- fit only what the capture determines --- and they
share neither a mechanism nor a failure mode. The angular mask also opens a memory result the anisotropic filter has no route
to, since masked coefficients can go unstored.}

\section{Normalising the threshold}
\label{sec:b2scale}

Equation~\eqref{eq:lmax} as written thresholds $\lambda_{\min}(A_k(\ell))$
against an absolute $\tau$, and that does not survive contact with a real
capture. Two separate scale dependencies are at fault, and they need separate
fixes.

\paragraph{View count.} $A_k$ is a \emph{sum} over views, so a primitive seen by
$100$ cameras has eigenvalues an order of magnitude larger than one seen by $10$,
before any question of conditioning arises. The normalised form
\begin{equation}
  \bar A_k \;=\; \frac{A_k}{\sum_n w_{nk}}
  \label{eq:gramnorm}
\end{equation}
removes it: $\bar A_k$ is a weighted \emph{average} of $y y^{\!\top}$ over the
directions that saw the primitive, so $\tau$ means the same thing for primitives
with different view counts. Without this, $\tau$ silently encodes a preference
for well-observed primitives that has nothing to do with whether their angular
sampling is well spread.

\paragraph{Angular spread.} Normalising is necessary, and on its own it is insufficient. Even on
$\bar A_k$, $\lambda_{\min}$ still scales with how widely the directions are
spread, so one $\tau$ cannot serve a $100$-camera orbit and a $10$-camera cone.
Section~\ref{sec:b2amend} measures this: at $\tau = 0.01$ the absolute form keeps
everything on a full orbit and nothing on either degraded protocol, and at
$\tau = 0.001$ it keeps $82\,\%$ on the arc. It is a cliff. The scale-free condition number
\begin{equation}
  \ell_{\max}(k) \;=\; \max\Big\{\, \ell \;:\;
  \frac{\lambda_{\min}\big(\bar A_k(\ell)\big)}{\lambda_{\max}\big(\bar A_k(\ell)\big)} > \tau \,\Big\}
  \label{eq:lmaxcond}
\end{equation}
behaves as the criterion should at a single threshold, and is what the
implementation uses. Equation~\eqref{eq:lmax} is retained above because it is
the form the argument is most naturally stated in, and because the amendment is
a result: it is a change to the method that only measurement could have
produced, and \S\ref{sec:b2amend} reports the measurement that forced it, with the
derivation amended in the open.

\paragraph{Status: implemented, measured, and refuted.} The angular mask is evaluated on all
eight Blender scenes and across the six capture protocols of
\S\ref{sec:stress} without a GPU, since Equation~\eqref{eq:gram} depends only on
primitive positions and training view directions; Table~\ref{tab:b2} adds the
rendered-quality and model-size columns.

The criterion does what this chapter says it does in the \emph{geometric}
sense: it retains everything on a full orbit, nothing on a pencil, and responds
gradually in between. Choosing \emph{which} coefficients to keep at a fixed budget is a separate
matter, and it is the claim the standard set below tests. Section~\ref{sec:btwo} reports the test and the criterion
fails it: against magnitude pruning restricted to the angular mask's own nested degrees, so that the structure and the budget are
the same and only the criterion differs, the angular mask loses by up to
$\btwoBlockedGapMax$\,dB. The derivation that follows is therefore presented as
a proposal that was carried through to measurement and did not survive it. It is
kept in full because the argument is sound and the conclusion is still
wrong, and the gap between the two is what a reader can learn from.

One caveat carries through to every statement about memory: masking coefficients
to zero does not shrink a fixed-width PLY, so \emph{no model-size saving is
claimed from these numbers}. Realising it needs variable-length SH storage,
which Chapter~\ref{ch:plan} schedules for Spring, and which, given
\S\ref{sec:btwo}, should be scheduled for a criterion that works.

\part{Evidence}
\chapter{Experimental design}
\label{ch:design}

\keybox{why the experiment is shaped like this}{%
Two ideas govern everything in this chapter.

The first is that a comparison is only as good as what it holds fixed. If you
compare two published methods by running each author's repository, the difference
between the two numbers contains the difference between two dataloaders, two
metric implementations and two perceptual-loss backbones, and you cannot tell
which part is the method. So every configuration here is the same checkout with
different flags, and the difference between two numbers is the flags.

The second is that a prediction made after seeing the data is not a prediction.
The acceptance levels, the decision gates, the thresholds and the compute budget
were all written down before any run, and are reproduced here unchanged. Two of
those gates could have ended the project --- if the resolution limit already binds
everywhere, or if real capture geometry turns out to be effectively isotropic,
then the method has nothing to act on. Both were designed to be cheap and to be
run first, before any method code existed. Chapter~\ref{ch:results} reports what
they returned.}

This chapter states the design: what is held fixed, what is varied, what would
count as success and what would count as failure. The commands that carry it
out, and the machinery that runs them, are Appendix~\ref{app:harness}.

\section{Compute envelope and what it forbids}
Free tier: Kaggle $2\times$T4, 16\,GB each, compute capability 7.5. 3DGS states
that 24\,GB is required for paper-quality training, so numerical identity with
the published tables is unavailable for the largest scenes. Three acceptance
levels are declared in advance.

\begin{table}[htbp]
  \centering
  \caption{Acceptance levels, declared before any run.}
  \label{tab:levels}
  \small
  \begin{tabular}{llp{0.42\linewidth}l}
    \toprule
    level & claim & criterion & where attainable \\
    \midrule
    L1 & numerical reproduction & mean PSNR within $0.20$\,dB, SSIM $0.005$, LPIPS(VGG) $0.010$ of published, at the paper's resolution, iterations and split & Blender synthetic \\
    L2 & phenomenon reproduction & the claimed effect reappears with correct sign, ordering and order of magnitude under a stated deviation & everywhere \\
    L3 & harness parity & own re-run baselines are the control for all method comparisons; published numbers cited, never used as the control & everywhere; governs the thesis tables \\
    \bottomrule
  \end{tabular}
\end{table}

L3 is standard practice: 3DGS itself marks its Mip-NeRF 360 row as copied and
reports its own re-run separately ($27.58$ against the copied $27.69$). An
independent prior on implementation identity is the ${\sim}0.6$\,dB spread
between faithful splatting implementations on the same benchmark --- three times
the L1 tolerance.

\section{Reproduction, and the single-codebase decision}
\label{sec:onecodebase}

\begin{figure}[htbp]
  \centering
  \includegraphics[width=0.99\linewidth]{figures/d8-two-arms.pdf}
  \caption{The construction every comparison in this thesis rests on. Four
  configurations, one checkout, and the only thing that differs between any two
  of them is the band-limit. Running each published method from its own
  repository would have been easier and would have meant that the difference
  between two numbers also contained the difference between two dataloaders, two
  metric implementations and two perceptual backbones, with no way to say which
  part was the method. Note that the published figures are used only as targets
  to check the reproduction against; every comparison reported in
  Chapter~\ref{ch:results} is between two of these columns.}
  \label{fig:twoarms}
\end{figure}
Both baselines run inside the Mip-Splatting codebase. The 3DGS arm is that
codebase with the 3D filter zeroed and \texttt{kernel\_size}~$=0.3$, which by
Proposition~\ref{prop:reduction} is 3DGS exactly.

This is forced as well as desirable. The multi-scale benchmark data is written
in mip-NeRF's multiscale format, and only Mip-Splatting's loader can read it;
the original 3DGS repository cannot open the dataset at all --- it falls
through every branch of the dispatch and dies. It also could not work in
principle: the multi-scale loader sets the field of view \emph{per frame},
whereas the Blender reader takes one global \texttt{camera\_angle\_x}.

The benefit is methodological: both arms share one dataloader, one metrics
implementation and one LPIPS backbone (VGG), so any difference between them is
the method rather than the plumbing. The diff that separates the arms is
enumerated in the kernel and asserted on every run, and a change outside it
aborts the session before any training starts.

\keybox{sharing a codebase shares more than the plumbing, and that is the trap}{%
``The 3D filter zeroed and \texttt{kernel\_size}~$=0.3$'' makes the arm 3DGS in
its \emph{3D} band-limit, exactly, and says nothing about its \emph{2D} one.
Both arms inherited Mip-Splatting's screen-space opacity compensation, which
3DGS does not have, so for the first three rungs the arm labelled 3DGS was
anti-aliasing better than 3DGS --- and the between-arm gap, the quantity this
whole construction exists to protect, carried a residual anti-aliasing term.
Section~\ref{sec:g2} is the diagnosis and the re-measurement. The lesson is not
that the single-codebase decision was wrong, because it is forced; it is that
``both arms share everything but the diff'' is a claim about what the diff
\emph{omits}, and omissions are exactly what a diff does not show.}

\paragraph{The arms, as run.} Arm A is Mip-Splatting unmodified. Arm B is the
same codebase with \texttt{filter\_3D} identically zero, \texttt{kernel\_size}
$=0.3$, \emph{and} the 2D opacity compensation disabled in the rasteriser --- all
three, which together are vanilla 3DGS's band-limit semantics in both domains.
The third was added after \S\ref{sec:g2} found it missing.

\section{The ladder}
\begin{table}[htbp]
  \centering
  \caption{The rungs, their budgets and their exit criteria. GPU-hours for R0
  and R1 are this project's measurements; the remainder are the estimates they
  replace.}
  \label{tab:ladder}
  \small
  \begin{tabular}{llp{0.44\linewidth}}
    \toprule
    rung & run & exit criterion \\
    \midrule
    R0 & smoke --- one scene, 7\,K, both arms, plus the twelve pre-flight assertions & toolchain fixed; estimates replaced by measurement \\
    R1 & Blender single-scale train $\rightarrow$ multi-scale test, 8 scenes $\times$ 2 arms & gate G2: $\tfrac18$-scale gap exceeds 9\,dB \\
    R2 & Blender multi-scale train and test & Mip-Splatting ahead at every scale; L1 within $0.20$\,dB \\
    R3 & real scenes that fit 16\,GB: 360 indoor $\times4$, T\&T $\times2$, DB $\times2$ & per-scene table; OOMs recorded, never worked around \\
    R5 & seed spread --- 3 seeds $\times$ 2 scenes $\times$ 2 arms & $\sigma$, without which parity is not a testable claim \\
    R4 & Mip-NeRF 360 outdoor at \texttt{images\_8} --- the only authorised deviation & L2 only, separate table, resolution in the caption \\
    \bottomrule
  \end{tabular}
\end{table}

\section{The stress suite}
\label{sec:stress}

\begin{figure}[htbp]
  \centering
  \includegraphics[width=0.99\linewidth]{figures/d4-protocols.pdf}
  \caption{The six capture protocols, drawn as what they are. The black dot is
  the object; the coloured dots are the training cameras. Only the training set
  changes --- the test set is the full orbit in every case, so the model is
  always asked about viewpoints it was not fitted from, and the six differ only
  in how much the capture knew. Two of them look alike on purpose: \emph{mixed
  focal} has the same camera positions as the full orbit and differs only in
  that half the images are downsampled (drawn hollow), which is a defect of
  camera \emph{settings} rather than of geometry; and \emph{grazing} keeps a wide
  span but selects only high-incidence views, which a top-down drawing cannot
  show. The bottom row is where the capture becomes genuinely
  ill-conditioned.}
  \label{fig:protocols}
\end{figure}
Equation~\eqref{eq:cap} predicts that the method's advantage is a function of
the angular span of the training cameras. Testing that requires captures of
varying conditioning, which are constructed as camera subsets of images already
downloaded --- no new data, no extra GPU time. Each protocol is versioned as a
JSON index list and carries its \emph{measured} span, because a subset built
with a nominal target of $6^\circ$ that comes out at $20^\circ$ is not the
capture the prediction was made for, and the prediction must then be re-evaluated
at the span that was actually built rather than the one that was asked for.

\input{generated/table-protocols}

\section{Three instruments, attached to the baseline runs}
\label{sec:instruments}
All read-only with respect to the optimiser, so the reproduction remains a
reproduction. Together they test the method's premise before the method exists.

\begin{table}[htbp]
  \centering
  \caption{The instruments and what each decides.}
  \label{tab:instrumentplan}
  \small
  \begin{tabular}{lp{0.34\linewidth}p{0.42\linewidth}}
    \toprule
    & measurement & what it decides \\
    \midrule
    I-1 & distribution of \texttt{filter\_3D} at end of training, per scene, against scene scale and distance to nearest camera & whether $s^2/\lambda_i$ ever exceeds $f_k^2$. If the floor binds almost everywhere, Proposition~\ref{prop:reduction} applies universally, the method collapses to the baseline, and the project must pivot. The cheapest possible falsification, and a by-product of a run already scheduled. \\
    I-2 & $\lambda_3/\lambda_1$ per primitive on trained, unmodified checkpoints, per capture protocol; pure PyTorch, no CUDA & whether the regime table of Equation~\eqref{eq:cap} holds on real scenes. Turns \S\ref{sec:anisotropy} from prediction into evidence before any method code is written. \\
    I-3 & per-scene mean and standard deviation across seeds, both baselines & the $\sigma$ that makes ``parity'' a quantity. Parity is then reported as within $k\sigma$ of the self-run baseline. \\
    \bottomrule
  \end{tabular}
\end{table}

One detail here has a real failure mode. A primitive that no training camera
sees has no distance to any camera, and the published code gives it the largest
distance in the scene rather than leaving it undefined. Those primitives must be
excluded from the instrument's histogram, or they pile up in its right tail and
read as though the resolution floor were binding hard --- which is precisely the
wrong answer at gate G3, and an encouraging-looking one.

\section{Metrics, baselines, statistics}
Quality is PSNR, SSIM and LPIPS (VGG) per scene and averaged, at test scales
$1\times$, $\tfrac12$, $\tfrac14$, $\tfrac18$. Cost is primitive count, model
MB, peak training VRAM, wall-clock and render fps --- same GPU, same run.
Diagnostics are the fraction of primitives where the estimation term exceeds the
floor, and the mean anisotropy of $\Sigma_{\mathrm{filt}}$; this is what makes a
parity result interpretable rather than mysterious. Statistics are three seeds
minimum, mean $\pm$ standard deviation, paired per-scene deltas.

\paragraph{One measurement the published code cannot make.} The multi-scale
test set holds all four resolutions mixed together, and the evaluation that ships
with both methods averages the whole mixture into a single number. But the papers
report four numbers, one per scale, and so does this thesis --- the entire
argument is about what happens as the test resolution leaves the training
resolution, which a single pooled average hides completely.

The four groups therefore have to be recovered. They can be, exactly and without
re-rendering, because the loader preserves the order in which the test views were
written and that order records which scale each view came from. Re-rendering at
each scale instead would \emph{resample} the images and answer a different
question, so it is not done. The recovery is checked rather than trusted: the
count-weighted mean of the four recovered groups must reproduce the pooled
number the shipped evaluation reports, and that is asserted on every scene of
every run.

\chapter{Results}
\label{ch:results}

\keybox{how to read this chapter}{%
The chapter runs in the order the work was done, which is also the order in which
each claim depends on the one before it.

First the reproduction: can this project produce the two published baselines at all? Nothing later means
anything without it. Then the two instruments, which measure the capture geometry itself and never a rendered image --- these are the
falsification gates, and they were run before the method existed. Then the method
on the standard benchmark, where the theory demands it change nothing. Then the
stress suite, where the theory says it should act, which is the decisive
experiment. Then the second method, which fails.

Three conventions hold throughout. Every number here is read from a single
append-only record and inserted mechanically; nothing is typed by hand,
and a quantity that has not been measured appears as a visible marker, never as a plausible figure. Published figures are cited as \emph{targets} and never
used as a control --- every comparison is against this project's own re-run
baseline. And where something went wrong, it is reported where it happened,
with the correction it forced; the engineering detail behind each is in
Appendix~\ref{app:defects} so that it does not interrupt the argument.}

\keybox{status}{%
Five reproduction rungs, two read-only instruments and both proposed methods
have been run, and everything below is measured. The protocol, acceptance
criteria, compute budget and decision gates were frozen \emph{before} any run,
so that a disappointing number later could not be rationalised into a success;
where a gate failed or a prediction did not hold, the chapter says so and keeps
the original wording. Every table and figure in this chapter is regenerated
mechanically from the measurement record; no number is typed into this document
by hand (Appendix~\ref{app:record}).}

\section{Pre-flight assertions}
\label{sec:preflight}

A reproduction that has not checked its own instrument is a reproduction of
nothing in particular. Twelve assertions therefore run before the first training
step, and the session aborts if any of them fails. They check that the GPU is
one the code can actually run on, that both arms load the same images, that the
perceptual metric uses the backbone the papers used, and that the two checkouts differ by exactly the intended change and
nothing else. Appendix~\ref{app:harness} lists all twelve with what each one
protects against. All twelve pass.
\section{Single-scale training}
\label{sec:table1}

\input{generated/table1-stmt}

\keybox{what reproduces, and to what level}{%
Over the $\stmtArmAScenes$ scenes the published mean covers, Mip-Splatting lands
within $\stmtArmAMaxDelta$\,dB of its published figure at every test scale.
That is inside the $\pm0.20$\,dB of Table~\ref{tab:levels}, so the Mip-Splatting arm carries an
\textbf{L1} claim on PSNR. The \textbf{gap} between the arms, on the
$\stmtMatchedScenes$ scenes both measured, is within $\stmtGapMaxDelta$\,dB of
the published gap at all four scales --- also L1, and it is the quantity the
two-arm construction exists to produce.

The 3DGS column is a different matter and claims nothing. \texttt{ship} does not
train on the vanilla arm (\S\ref{sec:shipfails}), so that column is a
$\stmtArmBScenes$-scene mean standing against an eight-scene published one. Its $\stmtArmBMaxDelta$\,dB is mostly the missing scene, and is reported as
such.}

\keybox{an artefact in an earlier version of this table}{%
Both \emph{ours} columns were computed on the matched seven scenes, and compared
against published figures that are means over eight. \texttt{ship} happens to be
the Mip-Splatting arm's weakest scene at full resolution, at $30.48$\,dB against a set mean of
$33.44$, so excluding it lifted the Mip-Splatting arm's mean by $0.42$\,dB, and that entire
lift was being reported as reproduction error. The Mip-Splatting arm read $0.50$\,dB from its
published figure and was written up as L2.

Two comparisons live in this table, and each needs its own denominator. A
column compared against a published eight-scene mean has to be an eight-scene
mean. A gap between two arms has to be over the scenes both arms have, or it is
partly the difference between two sets of objects. Computing both on the
matched set made the second correct and the first wrong, and the error was
invisible because it moved the number in the direction a reproduction error
would move it.}

\subsection*{The same result, as a picture}
\label{sec:qualitative}

A decibel is a summary, and a $10$\,dB gap in a table is a claim the reader has
to take on trust. Figure~\ref{fig:ladder} is the first row of
Table~\ref{tab:stmt} rendered instead of tabulated: one view of \texttt{lego},
trained once at full resolution, evaluated at each of the four test scales, on
both arms.

\begin{figure}[htbp]
  \centering
  \includegraphics[width=0.92\linewidth]{figures/fig11-qualitative-ladder.pdf}
  \caption{One view of \texttt{lego}, single-scale training, four test scales.
  Left: ground truth at that scale. Middle: 3DGS. Right: Mip-Splatting. The
  numbers are PSNR for the two images shown, recomputed from those pixels, which is why they differ from Table~\ref{tab:stmt}, a mean
  over $\stmtMatchedScenes$ scenes and every test view. Renders are magnified
  with nearest-neighbour interpolation so that a rendered pixel is a visible
  block: a $\tfrac18$-scale render is $100\times100$, and any smooth resampler
  would average away the aliasing the figure exists to show.}
  \label{fig:ladder}
\end{figure}

At the training scale the two are indistinguishable, and the numbers agree to
$0.07$\,dB. One octave down they have separated by $7$\,dB, and what separates
them is visible without being told where to look: the 3DGS reconstruction
\emph{thickens}. Treads merge into a band, the gap between the shovel arms
closes, the yellow spreads past the studs that bound it. By $\tfrac18$ the
machine is a blob with a red dot on it, and Mip-Splatting still resolves the
tracks.

That thickening is the mechanism of Chapter~\ref{ch:background} in visible
form. A primitive smaller than a pixel at the test scale is still splatted, at full
opacity, into the pixel that contains it, and many such primitives sum. The image gains energy it should not have, so
edges bloom outward and dark regions lift. The band-limit is what stops a
primitive from being smaller than the sampling can support, and the right-hand
column is what that buys.

\begin{figure}[htbp]
  \centering
  \includegraphics[width=0.92\linewidth]{figures/fig12-qualitative-scenes.pdf}
  \caption{The $\tfrac18$ row across four scenes, same construction as
  Figure~\ref{fig:ladder}. \texttt{materials} is the clearest case: the
  spheres lose their specular highlights to bloom and the dark ones lift
  towards grey. Every scene measured under both arms shows the same signature;
  these four are shown because they are visually distinct from one another. The
  remaining scenes are kept as per-scene contact sheets alongside the record, so that the selection can be audited instead of trusted.}
  \label{fig:qualscenes}
\end{figure}

\keybox{what these figures show}{%
They are the pixels the PSNR column was computed from. Each rung deletes all but
the first few rendered views once the metrics have consumed them, and what
survives is a matched set of predictions and ground truths per model; these are those files, carried straight out of the run that produced them. The ground-truth image is
byte-identical between the two arms at every index, and that is checked on every
pair, so the two predictions are of the same view.

They are also \emph{single views}, and a single view is not evidence of a mean.
The tables remain the claim; these show what the claim looks like.}

\input{generated/table-cost}

\section{The scale-collapse gate}
\label{sec:g2}

The scale-collapse gate was declared in advance as: 3DGS falls $\gtrsim13$\,dB from full to
$\tfrac18$; Mip-Splatting falls $\lesssim6$\,dB; the gap at $\tfrac18$ exceeds
$9$\,dB. On R1 it failed, on two of its three criteria: the 3DGS arm fell $9.09$\,dB
where $13$ was required, and the gap at $\tfrac18$ was $4.07$\,dB against a
published $10.98$. The protocol's instruction when the scale-collapse gate fails is to stop and
diagnose in a stated order. That took two attempts, and the first was wrong.

\subsection*{The protocol's three checks, and a fourth}

\begin{enumerate}[leftmargin=1.4em,itemsep=2pt]
  \item \textbf{The per-scale recovery.} Its count-weighted mean reproduces
        the shipped evaluation's own pooled figure to within $0.01$\,dB on every
        scene of every run --- asserted per scene, on every scene. Not the
        cause.
  \item \textbf{The 3DGS arm's \texttt{filter\_3D} at step 30\,000.} Read back from all
        eight saved point clouds it is identically zero, and non-zero on all
        eight of the Mip-Splatting arm's. The reduction is exact. Not the cause.
  \item \textbf{The white-background setting.} Passed to every job on both
        arms. Not the cause.
\end{enumerate}

\paragraph{The first diagnosis, and why it was wrong.} With those three
eliminated, the explanation offered was densification. The public Mip-Splatting repository
changed how it decides when to subdivide a primitive, roughly a year after the
paper was published: it adopted a criterion from Gaussian Opacity Fields, which
splits on a quantile over accumulated gradient magnitudes instead of on a fixed
gradient threshold. The original
3DGS rule is the fixed threshold, $\lVert\nabla\rVert \ge \tau_{\mathrm{pos}}$
alone. The argument was that a
better densification recovers most of what the missing band-limit costs, so the collapse the gate's thresholds were calibrated against could not occur in
this codebase, and the gate was therefore unmeetable here.

It was a plausible story, it was consistent with every number then available,
and it was \emph{wrong}. It also had the shape a wrong diagnosis usually has:
it explained the failure as unfixable, and it did so without a test that could
have refuted it.

\subsection*{The fourth check: the rasteriser on the 3DGS arm}
\label{sec:c1}

The three checks above say the gate failed. They do not say why, and the first
answer this project gave was wrong.

The 3DGS arm was built by zeroing Mip-Splatting's 3D smoothing filter, which makes the
world-space band-limit 3DGS's. What that overlooks is that Mip-Splatting has a
\emph{second} band-limit: a screen-space dilation carrying an opacity
compensation $\rho = \sqrt{\det\Sigma' / \det(\Sigma' + k I)}$. That one lives in
CUDA,
where zeroing a Python attribute does not reach it. The arm labelled 3DGS was therefore anti-aliasing better than 3DGS, at the very scales the gate was
scored on.

Isolating that term is what the opacity-compensation fix does. Appendix~\ref{app:impl} gives the switch and
what it touches; Appendix~\ref{app:defects} gives how it was found, including the
two further things the switch turned out to remove that nobody had asked it to.
\subsection*{The scale-collapse gate, re-scored}

Every 3DGS-arm row produced with the compensation active is superseded. They are
marked as such in the record and retained --- every table in this document
excludes them, and none of them has been deleted.
The Mip-Splatting arm is untouched and was not re-run: nothing about this defect involves it.

\input{generated/table-g2}

All three criteria pass, against thresholds that were set before any run and have
not been touched since. And the consequence reaches further than the gate: with
the arm rebuilt, Table~\ref{tab:stmt} reproduces \emph{both} baselines, where before it
reproduced one. 3DGS lands within $\stmtArmBMaxDelta$\,dB of its published figure at
every test scale and Mip-Splatting within $\stmtArmAMaxDelta$, and the gap column
--- the quantity the whole two-arm construction exists to measure --- reproduces
to within $\stmtGapMaxDelta$\,dB at all four, which is inside the L1 tolerance. This project did not have that before; it had the Mip-Splatting arm at L1 and a 3DGS arm several decibels adrift, for a
reason it had explained incorrectly.

\keybox{the cost of the rebuild}{%
It cost the R1 and R2 3DGS-arm columns and the 3DGS half of the seed spread --- three rungs of measurement, re-run, and it surfaced a numerical
fragility that had been hidden by the compensation (\S\ref{sec:shipfails}). It bought a 3DGS arm that is 3DGS in both of its band-limits, so the control every later comparison rests on is the method itself. The
between-arm comparison was the one thing the earlier diagnosis insisted was
safe, and it was the thing that was wrong.}

\subsection*{What the vanilla arm cost in robustness}
\label{sec:shipfails}

Removing the compensation cost something the protocol had not budgeted for. The
scene \texttt{ship} is much the largest in the suite, at $442\,235$ primitives
and $90\,\%$ of the card, and it does not train on the vanilla arm at all. It
has
been attempted nine times and has failed nine times, always with a CUDA illegal
memory access and never with an allocation error, at iterations spread from
$3\,000$ to $12\,100$. What varies is \emph{when}, not \emph{whether}.

Two things keep the exclusion narrow. \texttt{ship} trains
without incident on the \emph{multi-scale} protocol, so Table~\ref{tab:mtmt} is a
full eight-scene comparison and the failure is specific to single-scale training.
And it trains without incident on the Mip-Splatting arm, so the fragility belongs to the vanilla band-limit semantics.

It is excluded from the single-scale 3DGS-arm rows, and the protocol's rule for
a scene that will not run is to record it and move on. The instrumentation that
localised the fault, and the three explanations it ruled out, are in
Appendix~\ref{app:defects}.
\section{Multi-scale training}
\label{sec:table2}

\input{generated/table2-mtmt}

\paragraph{This column was two rasterisers, briefly.} The 3DGS arm was re-measured
after the opacity-compensation fix, and for a period each of its $32$ cells was the mean of the old measurement and the new one instead of the new one alone --- the rows the re-run
replaced were never retired, because the rule that retires them keyed on a name
the re-run shared. Nothing about the output looked wrong. The generator now
refuses to emit a table whose cells draw on more than one build, and names the
cell; the account is in Appendix~\ref{app:defects}.
\section{The densification explanation}
\label{sec:densification}

Section~\ref{sec:g2} settles the cause of the $\tfrac18$-scale shortfall: it was
the 2D opacity compensation, and removing it recovers the published shape on the
scene it was tested on. The densification observation stays true and turns out to be insufficient, so
what it implies needs keeping straight.

\keybox{the repository's densification}{%
\texttt{densify\_and\_prune} in this codebase splits on a quantile over
accumulated \emph{absolute} gradients. The quantity it thresholds is the
accumulated magnitude of the screen-space position gradient, and the threshold
is $Q = \mathrm{quantile}(\texttt{grads\_abs},\,1-\varrho)$, where $\varrho$ is
the fraction of primitives passing the ordinary gradient test. This is the Gaussian Opacity Fields rule, and it entered the public
Mip-Splatting repository about a year \emph{after} the paper's numbers were
produced. Original 3DGS splits on $\lVert\nabla\rVert \geq \tau_{\mathrm{pos}}$
alone. So the code every reproduction of Mip-Splatting now starts from does not
densify the way the published numbers did --- a real difference from both
published methods, and one both arms of this project share.}

Two independent reasons stop it from explaining the Mip-Splatting arm's distance
from its published figure. The first is that there is no longer much
distance to explain: on the eight scenes the published mean covers, the Mip-Splatting arm is
within $\stmtArmAMaxDelta$\,dB at every test scale. The $0.50$\,dB this section
was written to account for was the artefact of \S\ref{sec:table1} --- a
seven-scene mean standing against an eight-scene published one.

The second is that the hypothesis has now been tested directly, and it fails.

\input{generated/table-pregof}

The probe reverts exactly that one change and nothing else, $45$ lines across
two files, which puts the code back in the state Mip-Splatting's published table
was produced from. The Mip-Splatting arm was re-run on it, all
$\pregofScenes$ scenes, same data, same
seed, same iterations.

\keybox{the densification explanation}{%
At the pre-GOF commit the Mip-Splatting arm reproduces the published figures to
$\pregofMaxDelta$\,dB at every test scale --- effectively exactly, and well
inside L1. If GOF densification were the cause of a standing offset, the
post-GOF column would sit clearly above this one. It does not: the effect of
that change is at most $\gofMaxEffect$\,dB and it is not consistently signed,
$+0.07$ at full resolution and $-0.11$ at $\tfrac18$.

So the densification difference is real, and it is shared by both arms, but it
is \emph{not} why the Mip-Splatting arm sits where it sits. The explanation had never been
tested, and when it finally was, it turned out to be wrong. The sentence it
replaces read: ``a better densification plausibly accounts for standing above
the published level''. That sentence was doing the work of a measurement while
being an assertion, and the word \emph{plausibly} was carrying
all of it.}

The earlier version of this section also argued the point from the 3DGS arm's
$\cOneBaseFull$\,dB on \texttt{lego} at full scale, set against a published
eight-scene mean of $33.33$. That comparison is void: \texttt{lego} is one scene
and an easy one, and a single scene standing above an eight-scene mean says
nothing about either. It is the same mistake as the table above it, at a
different denominator, and finding one is what exposed the other.

Three consequences, revised:

\begin{enumerate}[leftmargin=1.4em,itemsep=3pt]
  \item \textbf{The Mip-Splatting arm carries an L1 claim on PSNR; the 3DGS arm carries none.} The Mip-Splatting arm is
        within $\stmtArmAMaxDelta$\,dB over the eight scenes the published mean
        covers. The 3DGS arm is missing \texttt{ship}, so its column is a
        $\stmtArmBScenes$-scene mean against an eight-scene published figure and
        is not a reproduction test at all. Only PSNR is compared: the published
        SSIM and LPIPS are not tabulated per scale in the source, so the L1
        claim is scoped to PSNR and stated that way.
  \item \textbf{The between-arm comparison is now what it was claimed to be.}
        Before the opacity-compensation fix the arms differed by the 2D compensation as
        well as by the 3D filter, so the gap column mixed the method with a
        residual anti-aliasing term. Both arms now share the densification, the
        dataloader, the metrics, the LPIPS backbone \emph{and} the screen-space
        filter semantics appropriate to each, and the difference between them is
        the band-limit.
  \item \textbf{The published $\tfrac18$ gap is reproducible after all.}
        Section~\ref{sec:g2} scores it. The earlier claim that it could not be
        reproduced in this codebase was a consequence of the defect, not of the
        codebase, and it is the clearest illustration in this project of why an
        explanation that makes a failure unfixable should be distrusted until it
        has been given a chance to fail.
\end{enumerate}

The shared-densification choice stands, and now for a better reason than before.
It was defended on the grounds that pinning the pre-GOF commit would move the
difference between the arms out of the filter and back into the plumbing, which
is still true. What has changed is that the cost of keeping GOF has now been measured instead of assumed, and it is $\gofMaxEffect$\,dB at worst.

\section{Run-to-run spread}
\label{sec:spread}
Parity needs a $\sigma$ before it is a testable claim. Repeats of an identical
configuration bound the harness's own nondeterminism, and repeats at different
seeds bound the variation the initial point cloud introduces.

\input{generated/table-seedspread}

Two different quantities are at stake here, and separating them turns out to
matter.

\begin{enumerate}[leftmargin=1.4em,itemsep=3pt]
  \item \textbf{Nondeterminism at a fixed seed} --- the same configuration run
        twice. CUDA atomics in the rasteriser accumulate in a non-deterministic
        order, so bit-identical inputs do not give bit-identical outputs. The
        largest per-scale difference observed is $\repeatSpread$\,dB.
  \item \textbf{Seed-to-seed variation} --- a different initial point cloud and a
        different view order. Three seeds on two scenes, both arms, all four
        scales, give a largest sample standard deviation of
        $\mathbf{\seedSigma}$\,dB, at \seedSigmaWhere.
\end{enumerate}

The second is $\seedVsRepeat\times$ the first, so the initialisation draw matters materially more than the atomics, which is
the reason the protocol marks R5 mandatory.

\keybox{the seed spread against the L1 tolerance}{%
L1 allows $\pm0.20$\,dB. A worst-case seed standard deviation of
$\seedSigma$\,dB is about half of that, so a single-seed L1 claim at an
unfavourable scale has a margin of roughly $2\sigma$ and no more. A three-seed mean has a standard error of $\seedSigma/\sqrt{3}$ and is
comfortable; a single run is adequate, with little room to spare. Both arms' reproductions in
\S\ref{sec:table1} are single-seed results, and several of their margins are of
the same order as this $\sigma$. They should be read as L1 attained, with the margin that implies and no more.

This figure is itself a consequence of \S\ref{sec:g2}: measured on the
compensated arm it was $0.139$\,dB, and on the vanilla one it is
$\seedSigmaArmB$. The seed spread of an arm is a property of that arm.}

The spread also differs by arm: the worst case over all four scales is
$\seedSigmaArmB$\,dB for 3DGS and $\seedSigmaArmA$\,dB for Mip-Splatting. The
arm without a band-limit is the less stable one under reseeding, which is
consistent with the mechanism this thesis proposes, though two scenes are far
too few to claim it as a finding.

\section{Real scenes}
\label{sec:real}

\input{generated/table-real}

Two results here, neither of which is a PSNR.

\paragraph{The compute envelope, resolved.} Section~\ref{ch:design} recorded 16\,GB
against a stated requirement of 24\,GB, and \S8.1 of the protocol listed
``whether \texttt{counter} and \texttt{drjohnson} fit 16\,GB'' as one of only two
facts that could not be settled without running. They can be settled now:
\texttt{counter} fits, at $14\,567$\,MB. Every scene measured peaks between
$14.5$ and $\realPeakVram$\,MB against a $15\,360$\,MB card --- under a gigabyte
of headroom, at up to $\realMaxGauss$ primitives. The envelope binds, and it holds: every scene ran at full settings, with nothing
reduced to make one fit.

\paragraph{The arms agree, which is the prediction.} On every one of these
scenes the two arms land within $0.2$\,dB of each other, and they do so in both
directions: 3DGS is ahead on \texttt{bonsai}, \texttt{counter} and
\texttt{kitchen}, Mip-Splatting on \texttt{room}. That is what Proposition~\ref{prop:reduction}
plus the floor-occupancy measurement of \S\ref{sec:instrumentresults} together require: these
are well-conditioned full-orbit captures, the Nyquist floor binds, and where it
binds the two filters are identical. It is a second confirmation of the instrument result, on real data, at the level
of rendered quality, and it is why no gain is available here for a method that
improves the filter only where the floor stops binding.

\paragraph{On comparison to published figures.} None is made. The published
Mip-NeRF 360 average is over nine scenes including the five outdoor ones, which
are excluded here by the memory envelope; a four-indoor subset average is a
different quantity and setting the two side by side is precisely the error
\S\ref{ch:discussion} records as having already occurred once in this project's
earlier write-up. These rows are L2 and are the self-run control for later work.

\section{Training cost and model size}
\label{sec:measuredcost}
The estimates in the protocol's \S7 were replaced by measurement at R0, on the
same free-tier hardware every later rung uses: $\rZeroItersPerSec$ iterations per
second on a Tesla T4, $\rZeroRenderFps$ rendered test views per second, and a
peak of $\rZeroPeakVram$\,MB of VRAM. The last of these matters for scope: 16\,GB
is nowhere near binding on the Blender scenes, so the compute envelope of
\S\ref{ch:design} constrains the \emph{real} scenes.

\section{Render speed}
\label{sec:fps}

Section~\ref{sec:implcost} argues that the anisotropic filter cannot cost anything at render time:
$\Lambda_k$ is recomputed on the filter's own training schedule and never inside
the render loop. That is an argument about code, and the first question a
graphics audience asks deserves a measurement.

\input{generated/table-fps}

The anisotropic filter renders at $\bOneFpsRatio\times$ Mip-Splatting --- $\bOneFpsSlowerPct\,\%$
slower, on a model with $\bOnePrimFewerPct\,\%$ fewer primitives, which is
within the variation two training runs of the same configuration produce in
primitive count alone. The filter's own construction costs $\bOneFilterSetup$\,s,
paid once before the loop begins, and $\Lambda_k$
is never touched inside it. So the claim of \S\ref{sec:implcost} rests on a measurement.

\paragraph{The missing 3DGS row.} The benchmark mounts each arm's own
trained model and scores whichever it finds. The vanilla 3DGS model was to come from the C2 session, which ended in error with six of eight scenes complete
(\S\ref{sec:budget}); the mount did not carry a \texttt{lego} model the
benchmark could discover, so that row is absent, and it has not been estimated. Nothing
about the comparison this section defends depends on it: the quantity at issue is
the anisotropic filter against \emph{its own baseline}, and both of those are here, measured in one
session on one GPU.

\keybox{what the \texttt{render\_fps} column measures}{%
It is views divided by the wall time of a whole rendering pass,
so it carries the process start, the scene load, the PLY parse, PNG encoding and
the disk writes, which is why it reads $8$--$10$ against a published
$180$--$300$, and why quoting it as a rendering result would be a serious
misreading of this project's own data. It is retained in Table~\ref{tab:cost}
because it is what that rung logged and deleting a logged number is not an
option; it is not used as a rendering claim anywhere. The measurement that
answers the question is Table~\ref{tab:fps}, and what it defends is the
\emph{ratio}, since every cost shared by the three arms has been excluded from
all of them equally.}

\section{Floor occupancy and conditioning}
\label{sec:instrumentresults}

Both instruments are functions of the primitive positions and the capture
geometry alone, so neither needs the images, the rasteriser, or a training run
of its own. Both were therefore run over every trained scene, and over the five
capture protocols of \S\ref{sec:stress}, at a total cost of two minutes.

\input{generated/table-instruments}

\begin{figure}[htbp]
  \centering
  \includegraphics[width=0.86\linewidth]{figures/fig8-floor-binding.pdf}
  \caption{The mechanism, measured. The anisotropic filter can only differ from Mip-Splatting where
  the estimation term $s^2/\lambda_i$ exceeds the Nyquist floor $f_k^2$; below
  the dashed line Proposition~\ref{prop:reduction} makes them the same filter to
  numerical precision. On the standard full orbit the floor dominates by a factor
  of three, and it still dominates under grazing and mixed-focal capture. Only
  the arc and the cone cross the line, which is the entire content of the
  claim that this method is a capture-conditioning method. The percentages under each protocol are the fraction of
  primitives above the floor in the worst-constrained direction and in all
  directions.}
  \label{fig:floorbinding}
\end{figure}

\paragraph{Two defects sit behind these numbers.} The instruments and the
training kernels were, for a period, cloning different branches, so a protocol
named \texttt{cone} meant one capture to the instrument and another to the run
that rendered it; and the constant $0.2$ of Equation~\eqref{eq:mipfilter} was carried
on one side of a comparison and not the other, which made the estimation term
appear five times larger than it is. Both are reported in
Appendix~\ref{app:defects}; the numbers in this section are from after both were
fixed, and the agreement between the instrument and the filter's own diagnostic
is the check that confirms it.
\subsection*{The floor-occupancy instrument}

The estimation term $0.2/\lambda_i$ and Mip-Splatting's floor $f_k^2$ are the two
sides of Proposition~\ref{prop:reduction}'s hypothesis, and the ratio between
them decides whether the method can act at all. Measured across the six
protocols, it rises monotonically as the capture closes: $0.16$ on the full
orbit, $0.21$ under mixed focal, $0.24$ under grazing, $0.48$ on the arc, and
past unity only on the cone ($1.35$) and the pencil ($5.52$).

On the standard benchmark the floor dominates by a factor of six, and the
estimation term exceeds it on \textbf{zero} primitives, in \emph{any} direction,
of \emph{every} scene. By Proposition~\ref{prop:reduction} that makes the anisotropic filter \emph{be} Mip-Splatting on this benchmark, exactly:

\keybox{why the standard benchmark offers no gain}{%
The eight Blender scenes are 100-camera full orbits, the best-conditioned capture
in the suite. A method evaluated only on this benchmark could not show an
improvement, and if one appeared it would be a bug. This is the cheapest
falsification the project had, it was run before a line of method code existed,
and Table~\ref{tab:method} is the confirmation at the level of rendered quality:
$\methodParityMax$\,dB over $\methodParityScenes$ scenes and four test scales.}

\subsection*{Agreement between the two diagnostics}

The floor-occupancy instrument is computed off a trained point cloud by a NumPy instrument; the filter
prints its own occupancy during training, from the model it is actually
filtering. They are independent implementations of the same comparison, and
Table~\ref{tab:stress}'s \emph{above floor} column can be read against
Table~\ref{tab:instruments}'s \emph{worst dir} column directly.

They agree on five protocols of the six. Both read $0\,\%$ on the full orbit,
on grazing, on mixed focal and on the arc, and both put the pencil above the
floor on essentially everything, at $100\,\%$ against the filter's $93\,\%$.

The one disagreement is the cone, where the instrument reads $100\,\%$ and the
filter $0\,\%$. That looks alarming until you notice that the ratio of
estimation term to floor is $1.35$ there, which is within $35\,\%$ of unity. A
verdict balanced that finely ought to be sensitive to what it is computed on, and the two are not computed on the same thing. The instrument
takes a model trained on the \emph{full orbit} and asks what the cone's ten
cameras would have been able to determine about it. The filter takes the model
those ten cameras actually produced, which has different primitives sitting at
different depths and therefore a different floor. Each is right about its own
question. The one that decides what the filter does at training time is the filter's, so that is the one the stress table reports.

\keybox{what the constant $0.2$ bought}{%
Before the constant was corrected (\S\ref{sec:m2defect}) these two columns
disagreed on the arc as well as the cone, and by a margin no point-cloud
difference could explain. Two implementations of one comparison agreeing on five
of six protocols, and differing only where the quantity sits within a third of
its decision threshold, is a much stronger statement than either makes alone.}

\subsection*{The conditioning instrument across the stress suite}

The refutation condition at the floor-binding gate is that the floor exceeds the estimation
term on nearly every primitive in \emph{every} scene, which would make Proposition~\ref{prop:reduction} apply universally and collapse
the method to its baseline everywhere. The measurement rules that out. Degrading
the capture moves the balance as Equation~\eqref{eq:cap} requires, and the measured anisotropy tracks the
prediction across a factor of eight:

\begin{itemize}[leftmargin=1.4em,itemsep=2pt]
  \item \textbf{Full orbit} ($179^\circ$): $\instFullAniso\times$ measured
        against $\instFullCap\times$ predicted.
  \item \textbf{Grazing} ($168^\circ$): $\instGrazingAniso\times$ against
        $\instGrazingCap\times$ --- the largest relative residual in the table,
        and the protocol built to maximise obliquity, which is what
        \S\ref{sec:offaxis} predicts a parallax-only model will miss.
  \item \textbf{One-sided arc} ($94^\circ$): $\instArcAniso\times$ against
        $\instArcCap\times$.
  \item \textbf{Low-parallax cone} ($51^\circ$): $4.01\times$ against
        $3.16\times$, and the first protocol where the estimation term exceeds
        the floor along the worst-constrained direction.
  \item \textbf{Narrow pencil} ($20^\circ$): $9.42\times$ against $8.22\times$,
        and the estimation term exceeding the floor by a median factor of five.
\end{itemize}

\keybox{the cap form as the prediction}{%
Every measured anisotropy comes out above the two-view prediction of
Equation~\eqref{eq:anisotropy}, and by roughly the same factor each time. The
fact that the factor is roughly constant is the clue: this is one discrepancy observed five times, and \S\ref{sec:capprediction} says
what it is. A capture protocol is a whole distribution of cameras spread over an angular window, where the two-view form places them at its two ends. For such a distribution the
correct $\lambda_3/\lambda_1$ is asymptotically half the two-view value, which
means the two-view form understates $\sigma_D/\sigma_L$ by a factor of
$\sqrt{2}$.

Scored against the cap form instead, at each protocol's own measured span, the
$\capAgreeN$ protocols that form is built for agree to
$\capAgreeMin$--$\capAgreeMax\,\%$. The two-view expression, scored the same way
on the same measurements, falls short by as much as $\twoViewLowMax\,\%$.

The sixth protocol, \emph{grazing}, sits at $\capAgreeGrazing\,\%$, which is the expected result. The cap form
describes parallax and nothing else, and grazing is the protocol built
specifically to test obliquity at close to full parallax. It is the row the model
does not claim to cover.

One detail supports the same reading. The residual is one-signed: measurement
sits above prediction every time, never below. That is what
\S\ref{sec:offaxis} would lead you to expect, because perspective contributes
anisotropy of its own that a cap of ray directions does not contain.}

So the floor-binding gate and the anisotropy gate both pass, and the regime
table of \S\ref{sec:anisotropy} rests on measurement --- before a line of method
code exists. This is the cheapest
falsification the project had available, and it did not fire.

\paragraph{The mixed-focal prediction.} Section~\ref{sec:defects} argues that
mixing cameras is the third defect of the scalar filter, and \S\ref{sec:stress}
expected the mixed-focal protocol to expose it. The measurement shows otherwise. The anisotropy moves
from $\instFullAniso\times$ to $\instMixedAniso\times$, which is nothing, and the
median ratio of estimation term to floor from $0.16$ to $0.21$ --- a real shift,
and still a factor of five below unity. Section~\ref{sec:table3} confirms it at
the level of rendered quality on a protocol that genuinely downsamples half the
cameras: the floor binds on every primitive and the two arms differ by
$+0.017$\,dB. The honest reading is that the mixed-camera defect is a statement about the \emph{value} of $f_k$, and that a $2$--$4\times$ focal spread stays too narrow to make it bind. Either the
protocol needs a wider spread or the defect is real but small; this measurement
does not distinguish those, and it is not reported as though it did.

\section{The standard benchmark}
\label{sec:table2m}

\input{generated/table-method}

\begin{figure}[htbp]
  \centering
  \includegraphics[width=\linewidth]{figures/fig9-parity.pdf}
  \caption{Proposition~\ref{prop:reduction} as a measurement. \emph{Left:} the
  three methods over the scenes measured under all of them --- the anisotropic filter's curve is
  drawn dashed \emph{on top of} Mip-Splatting's, and is invisible beneath it
  because that is what the proposition requires. \emph{Right:} the same thing as
  a difference, against the harness's own run-to-run spread. Every scale sits
  inside $\pm\sigma$. The 3DGS curve is there to set the scale: the difference
  the band-limit makes is several decibels, and the difference between
  \emph{which} band-limit is hundredths of one.}
  \label{fig:parity}
\end{figure}

\keybox{parity, predicted and recorded first}{%
The floor-occupancy instrument measured the Nyquist floor binding on $100\,\%$ of primitives on these
captures (\S\ref{sec:instrumentresults}), and Proposition~\ref{prop:reduction}
then makes the anisotropic filter and Mip-Splatting \emph{identical}. So this table is where the method has to match its baseline exactly. A gain here
would indicate a bug, and the exact-reduction unit test is run
inside the same session that produces these numbers, aborting the run if it
fails.}

\subsection*{What parity costs}
\label{sec:methodcost}

Parity is only half a result. A method that matches its baseline exactly and
costs nothing to run is free; one that matches it exactly and costs three times
the memory has to justify itself somewhere else. Until the paired session of
\S\ref{sec:table2m} there was no run in which both methods trained the same
scenes on the same card in the same session, so cost could only have been
compared across runs --- the comparison this project has had to correct three
times.

\input{generated/table-method-cost}

\keybox{matched primitive count, higher memory}{%
The anisotropic filter densifies to within $0.2\,\%$ of Mip-Splatting's primitive count, which is
Proposition~\ref{prop:reduction} arriving from a third direction: where the
floor dominates the filter is the same filter, so the optimisation follows the
same trajectory and produces the same model. The model on disk is the same size,
and rendering is unaffected --- the filter is applied once at construction (\S\ref{sec:fps}).

Training is another matter. The anisotropic filter costs $2.8\times$ the peak VRAM and $1.32\times$
the wall-clock time, for a result that is by construction identical. The
accumulation of $\Lambda_k = \sum_n w_{nk}(J_nW_n)^{\!\top}\Sigma_{\mathrm{pix}}^{-1}(J_nW_n)$
is a $3\times3$ symmetric matrix per primitive over every training camera, and
the peak does not track the primitive count. \texttt{ficus} at $187$k primitives
and \texttt{mic} at $406$k both peak near $13$\,GB, which is what a large transient in that accumulation looks like.

That is the cost line. A practitioner reading
this table on a well-conditioned capture should not use the anisotropic filter: it is the same
model, slower, in more memory. The case for it is confined to the regime of
\S\ref{sec:table3}, and \S\ref{ch:discussion} does not argue otherwise.}

\keybox{the missing scene}{%
\texttt{ship} was scheduled in this session and only Mip-Splatting completed it.
The cause was a scheduling race: both jobs started on \texttt{ship}
simultaneously on the two GPUs, the first pair in the queue to share a scene,
and raced on the point cloud the loader writes into the source directory --- the anisotropic filter
aborted in \texttt{create\_from\_pcd} with a null point cloud, four seconds in.
That is a harness defect, since fixed by ordering the job list so that adjacent
jobs never share a scene.

It is not re-run, and the reason is this table.
Mip-Splatting alone peaks at $14\,459$\,MB on \texttt{ship}, against a
$15\,360$\,MB card. At the $2.8\times$ measured here, the anisotropic filter on \texttt{ship} does
not fit, and no ordering of the queue changes that. The paired parity figure is
therefore over $\methodParityScenes$ scenes and says so.}

\paragraph{Three bugs on the way here.} Requiring parity \emph{exactly} turned the reduction theorem into a test, and the test
failed three times before it passed --- a clamp that was doing real work, a
determinant computed in a way that lost precision at realistic scales, and an
isotropic branch that re-derived a quantity it should have taken. Each produced
plausible numbers and the wrong shape. The diagnosis of each, and the unit test
that now covers the range where they hid, is Appendix~\ref{app:defects}.
\section{The stress suite}
\label{sec:table3}

Where the claim lives. Training cameras are restricted to each protocol's subset
while the test set is left whole, so the only thing that varies between rows is
the geometry the reconstruction had available. Both methods are trained inside
one session on the same data with the same seed, so the delta is paired within scene and test scale, never taken between two
means.

\input{generated/table-stress}

\keybox{how to read this table}{%
Proposition~\ref{prop:reduction} is a constraint on this table, which is a stronger thing than a prediction about
it. Where the estimation term does not exceed the Nyquist floor in
any direction, $\Sigma_{\mathrm{filt}} = f_k^2 I$ identically and the anisotropic filter \emph{is} Mip-Splatting: the same filter, taking the same code path.
On such a row $\Delta$ is a measurement of CUDA atomics and nothing else, and a non-zero value there would be evidence of a bug. So the
diagnostic column decides whether the PSNR column is an experiment at all, which
is why it is printed by the filter during training and carried here beside the result instead of reconstructed afterwards.}

\subsection*{What the table says}

Read down the \emph{above floor} column and the $\Delta$ column together and the
result is a single sentence: \textbf{the method does nothing on five protocols
and helps on the sixth, and the five where it does nothing are the five where Proposition~\ref{prop:reduction} says it must.}

On the full orbit, the grazing orbit, the mixed-focal capture, the one-sided arc
and the count-based cone, the estimation term exceeds the Nyquist floor on
\emph{zero} primitives. $\Sigma_{\mathrm{filt}} = f_k^2 I$ identically, the mean
anisotropy is $1.00\times$ to two decimals, and the two arms take the same code
path. The measured differences over those $\stressFlatCount$ rows run from
$\stressFlatMin$ to $\stressFlatMax$\,dB --- every one inside the harness's own
run-to-run spread, and none
exceeding $0.8\sigma$. That is the proposition holding, on $\stressFlatCount$ independently degraded
captures, to the precision of the instrument.

The narrow pencil is the first capture in this project on which the filter
actually engages. The estimation term exceeds the floor on $93\,\%$ of
primitives, and the filter comes out anisotropic, at $1.47\times$.
There, and only there, the anisotropic filter is ahead. The margin is $\stressBindDelta$\,dB on the
mean, which is $\stressBindSigma$ times the run-to-run spread. That mean is made of
eight paired cells, two scenes by four test
scales, and the anisotropic filter leads in every one of them, with per-cell gains running from
$+0.163$ to $+0.217$\,dB. A mean can come out positive by luck. A sign that
holds across all eight pairings cannot.

\begin{figure}[htbp]
  \centering
  \includegraphics[width=0.92\linewidth]{figures/fig13-qualitative-pencil.pdf}
  \caption{The operating point where the anisotropic filter acts, rendered. Three training cameras
  spanning $29^\circ$; the test set is the full orbit, so most of what is
  evaluated was never observed. Numbers are PSNR for the images shown. The anisotropic filter leads on both scenes, and on all four test scales of
  each, and the lead is too small to see.}
  \label{fig:qualpencil}
\end{figure}

\keybox{how to read these renders}{%
Both reconstructions are ruined. The chair is a smear with floaters where the
unobserved sides should be, the \texttt{lego} bulldozer is barely identifiable,
and nothing in either image would be accepted as output. The anisotropic filter's advantage is real,
signed consistently, and $\stressBindSigma\times$ the run-to-run spread --- and
it is also a fifth of a decibel on a $14$\,dB reconstruction, which is invisible.

Both statements are true, and the figure is here so that the first is not read
without the second. What Table~\ref{tab:stress} establishes is that the
\emph{mechanism} is the one the derivation names: the gain appears exactly where
the estimation term overtakes the floor and nowhere else. What it does not
establish, and what this figure makes plain, is that the method is useful at
this operating point. A capture this degraded needs more cameras.}

\keybox{the predicted ordering}{%
Section~\ref{sec:anisotropy} predicts \emph{where} this method helps: it stays
inert wherever the Nyquist floor dominates, and acts more strongly as the
capture's angular span closes. A method that helped uniformly would contradict
its own derivation, and \S\ref{ch:discussion} recorded that in advance as a reason for suspicion. What the table shows is the
predicted step: five rows inside the noise, one at four times it, and the step
falls exactly at the span where the estimation term overtakes the floor. The conditioning column rises monotonically alongside it, through $1.17$,
$1.22$, $1.48$, $2.47$, $4.11$ and $6.17\times$, and is measured on the same captures the PSNR is, at the cost
\S\ref{sec:branchdefect} records. The
crossover is between the count-based cone at $54^\circ$ and the pencil at
$29^\circ$, and the suite contains both only because the two protocol
definitions of \S\ref{sec:branchdefect} disagreed and both were kept.}

\subsection*{The limits of this result}

\begin{enumerate}[leftmargin=1.4em,itemsep=3pt]
  \item \textbf{The regime is narrow, and the absolute numbers say so.} A
        three-camera capture reconstructs badly: both arms sit near $14$\,dB on
        the pencil, against $34$ on the full orbit. The gain is real and it is
        measured on a reconstruction nobody would ship. What this establishes is
        that the mechanism is the one claimed, namely that the filter acts when and only when
        the estimation term binds. It does not establish that the method is
        useful at this operating point.
  \item \textbf{Two scenes, one seed.} The eight consistent pairings are eight
        cells of one run, not eight independent experiments. The seed spread of
        \S\ref{sec:spread} bounds the noise on a cell; it does not bound the
        scene-to-scene variation of the effect, and two scenes cannot.
  \item \textbf{The count-based cone is a null result, not a missing one.} At
        $54^\circ$ the floor still binds everywhere and the method is provably
        inert. That row is evidence about where the crossover is, and it is the
        reason the crossover can be stated as a measurement instead of an extrapolation from
        the pencil alone.
\end{enumerate}

\begin{figure}[htbp]
  \centering
  \includegraphics[width=0.9\linewidth]{figures/fig7-stress-delta.pdf}
  \caption{The central empirical claim, in one panel. Paired PSNR difference
  between the anisotropic filter and Mip-Splatting per capture protocol, with the harness's own
  run-to-run spread shaded at $\pm3\sigma$; a bar inside the shaded band is not a
  result. Protocols are ordered by conditioning, best at the top, and each
  carries the fraction of primitives on which the estimation term exceeded the Nyquist floor, the quantity that determines whether the two arms
  were running the same filter at all.}
  \label{fig:stress}
\end{figure}

\section{The angular mask}
\label{sec:btwo}
\label{sec:table4}

Equation~\eqref{eq:gram} depends on primitive positions and training view
directions alone, so the criterion is measurable from a trained model with no
retraining and no GPU. It was evaluated on all eight Blender scenes and then, for
the same reason the anisotropic filter had to be, across all six capture protocols.

\keybox{the angular mask on the standard benchmark}{%
Every primitive in all eight scenes is seen from $100$ well-spread directions,
and the criterion keeps degree 3 everywhere: mean $\ell_{\max} = 3.00$,
$100\,\%$ of non-DC coefficients retained. That is the correct answer, and it is
the same shape as the floor-occupancy instrument's result for the anisotropic filter --- on a well-conditioned capture the
proposed criterion reduces to the baseline. Neither contribution can be
demonstrated on a full orbit, and this has now been measured for both, where before it was argued for one.}

\input{generated/table-b2-protocols}

\begin{figure}[htbp]
  \centering
  \includegraphics[width=0.82\linewidth]{figures/fig10-b2-retention.pdf}
  \caption{The angular mask across the capture protocols. The same shape as
  Figure~\ref{fig:floorbinding}, in the angular domain instead of the spatial one: on a well-conditioned capture the criterion keeps every coefficient and
  the method reduces to its baseline, and it removes coefficients only where the
  views cannot support them. A criterion that pruned on the full orbit would be wrong.}
  \label{fig:b2retention}
\end{figure}

Under degraded capture the criterion does what \S\ref{ch:method2} claims: a
primitive seen from a narrow cone cannot support degree 3, and is not given it.
The retained fraction falls monotonically with the conditioning, from $100\,\%$
on the full orbit to $16.6\,\%$ in the cone.

\subsection*{An amendment to Equation~\eqref{eq:lmax}}
\label{sec:b2amend}

As \S\ref{ch:method2} writes it, the criterion thresholds
$\lambda_{\min}(A_k(\ell))$ against an absolute $\tau$. Measured, that does not
work across regimes. $\lambda_{\min}$ carries units and scales with the angular
spread of the views, so one $\tau$ cannot mean the same thing on a 100-camera
orbit and a 10-camera cone. On \texttt{lego}:

\begin{table}[htbp]
  \centering
  \caption{The absolute criterion has no $\tau$ that responds gradually. Fraction
  of non-DC coefficients retained.}
  \label{tab:b2abs}
  \small
  \begin{tabular}{lrrr}
    \toprule
    protocol & $\lambda_{\min} > 0.01$ & $\lambda_{\min} > 0.001$ & $\lambda_{\min}/\lambda_{\max} > 0.01$ \\
    \midrule
    full orbit        & $100.0\,\%$ & $100.0\,\%$ & $100.0\,\%$ \\
    grazing           & $94.3\,\%$  & $100.0\,\%$ & $100.0\,\%$ \\
    one-sided arc     & $0.0\,\%$   & $82.4\,\%$  & $44.6\,\%$ \\
    low-parallax cone & $0.0\,\%$   & $20.0\,\%$  & $13.3\,\%$ \\
    \bottomrule
  \end{tabular}
\end{table}

At $\tau = 0.01$ the absolute form keeps everything on a full orbit and
\emph{nothing} on either degraded protocol; at $\tau = 0.001$ it keeps $82\,\%$
on the arc, which is nearly everything. It is a cliff.
The condition number $\lambda_{\min}/\lambda_{\max}$ is scale-free and behaves as
the criterion should at a single threshold, which is why it is the default in the
implementation and why Equation~\eqref{eq:lmax} should be restated in that form.
This is a change to the method that only measurement could have produced, and it is recorded as such instead of folded silently into the derivation.

\subsection*{The full-orbit comparison}

\input{generated/table-b2}

Every cell in that table is identical across the three rows, and that is the result. On the standard benchmark the
criterion retains $100\,\%$ of non-DC coefficients, so the masked model \emph{is} the control, the same file scored through the same
render and metrics path, and magnitude pruning matched to the same fraction
prunes nothing either. The $\tau$ sweep says the same thing from another direction:
between $\tau = 0.003$ and $\tau = 0.1$, a factor of thirty, nothing on a full
orbit changes.

\keybox{the comparison this table cannot make}{%
Table~\ref{tab:b2} exists to ask whether identifiability masking beats magnitude
pruning \emph{at matched model size}. That question has no content when the
criterion prunes nothing: both rows keep everything, so both score identically,
and the table cannot distinguish a good criterion from a bad one. Reporting the
three equal rows as evidence that the angular mask matches magnitude pruning would be the error this chapter warns about elsewhere --- a condition tested where it was
constrained to be satisfied. So the comparison was run again on a capture where
the criterion does mask.}

\subsection*{The comparison on a capture where the angular mask prunes}
\label{sec:b2loses}

Restricting the cameras the \emph{criterion} reads to the low-parallax cone, and
leaving the test set whole, the angular mask retains $13$--$20\,\%$ of non-DC coefficients.
Magnitude pruning is then matched to that exact fraction, so the two models are
the same size and differ only in which coefficients survive.

\input{generated/table-b2-matched}

\keybox{the angular mask against magnitude pruning}{%
On every configuration where both methods actually prune, magnitude pruning
wins. Against the free control it wins by $0.60$ and $0.82$\,dB on \texttt{lego}
and by $1.82$\,dB on \texttt{chair}; against the \emph{fair} control below, by
$0.21$ and $0.42$\,dB on \texttt{lego} and by $0.09$\,dB on \texttt{chair}.
Chapter~\ref{ch:method2} states the standard this was to be held to: ``a criterion that cannot beat
magnitude pruning at matched size is not worth the machinery''. By that standard
the criterion fails, on the fair comparison
as well as the unfair one. The row at $0\,\%$ is both methods discarding
everything and agreeing by definition; it is not evidence either way.}

\paragraph{Why, and what survives.} Identifiability and importance are different
quantities, and the metric rewards the second. A coefficient can be badly
conditioned by the capture and still carry most of the fitted energy for that
primitive; zeroing it costs more PSNR than zeroing a well-determined coefficient
that happens to be tiny. Magnitude pruning optimises directly for what PSNR measures; the angular mask
optimises for something else and is then scored on PSNR.

\paragraph{The confound in the first comparison.} A second asymmetry sits in the
free control, and it is a large one: the angular mask masks whole
\emph{degrees}, removing coefficients in blocks of $2\ell+1$ nested from the top
down, while free magnitude pruning ranks individual coefficients and keeps the
largest anywhere. At the same budget that is strictly more freedom, so a loss against it is evidence about the \emph{structure} of the mask, when
the criterion is what is on trial. The first version of this
section noted the confound and left it, which was the right thing to notice and
the wrong place to stop.

The control that separates them is magnitude pruning restricted to the angular mask's own
nested degrees: same structure, same budget, and the only remaining difference
is what the two rank by. It is the \emph{blocked} column of
Table~\ref{tab:btwomatched}, and it moves the answer a long way. On
\btwoStructureScene, where the free control had beaten the angular mask by
$\btwoFreeGap$\,dB, the fair one beats it by $\btwoBlockedGapMin$ --- inside the
seed spread of $\seedSigma$\,dB, so on that scene the criterion and magnitude pruning are indistinguishable. $\btwoStructureShare\,\%$ of what looked
like a verdict on identifiability was a verdict on block structure.

What does not change is the direction. On \texttt{lego} the fair control still
wins at both budgets where anything is pruned, by up to
$\btwoBlockedGapMax$\,dB --- four times the seed spread. So the honest statement is narrower and
firmer than the one it replaces: at matched structure and matched budget,
ranking spherical-harmonic coefficients by how well the capture determines them
is worse than ranking them by how large they are --- by a fifth of a decibel
instead of two, decisive on one scene of the two and inside the noise on the
other.

What survives is narrower than what was claimed. The criterion still does what \S\ref{sec:b2scale} says it does: it responds
gradually to the conditioning, retaining $100\,\%$, $62\,\%$ and $25\,\%$ as the
capture degrades. That is a statement about the \emph{geometry}, measured and settled. What is not
supported is the step from there to rendered quality at matched size. On the
evidence here, identifiability is a good description of what the capture
determines and a poor rule for deciding what to keep.

\paragraph{Model size on disk.} Model size holds still, and has to: masking coefficients to zero leaves a fixed-width PLY the same size on
disk. Realising the saving needs variable-length SH storage, which
Chapter~\ref{ch:plan} schedules for Spring. The $74.8$\,MB in all three rows is
the same file written three times.

\section{What is established}

The reproduction rungs establish a control. Both baselines now exist as this project's own
measurements, on known hardware, from pinned commits, with the arm separation verified on the saved models. Every later comparison is against these numbers, which is
the L3 discipline of Table~\ref{tab:levels}.

About the proposed method they establish four things and refute one. The anisotropic filter is
implemented and reduces to its baseline exactly where the theory says it must,
verified to $1.4\times10^{-21}$ in the unit test and to $\methodParityMax$\,dB, comfortably below the seed spread, over
$\methodParityScenes$ scenes and four test scales in paired training runs. The conditioning both methods depend on
varies across capture protocols as Equation~\eqref{eq:cap} predicts.
The stress suite has been run, and the prediction recorded before it holds: on
the $\stressFlatCount$ protocols where the Nyquist floor still dominates the
difference is $\stressFlatMin$ to $\stressFlatMax$\,dB and the filter is
isotropic to two decimals, and on the one where the estimation term binds the anisotropic filter
leads by $\stressBindDelta$\,dB, on the mean and on every paired cell. And the cost of that has now been measured, where before it was assumed: $2.8\times$ peak VRAM and
$1.32\times$ training time for a model with the same primitive count
(\S\ref{sec:methodcost}).

What is refuted is the angular mask. At matched structure and matched budget it loses to
magnitude pruning, by $\btwoBlockedGapMax$\,dB at worst, and the criterion is
reported as failing the standard Chapter~\ref{ch:method2} set for it
(\S\ref{sec:btwo}).

What is \emph{not} established is that the anisotropic filter is useful. The regime where it helps
is one in which both methods produce reconstructions nobody would ship
(Figure~\ref{fig:qualpencil}), and on well-conditioned captures it is the same
model for three times the memory. Chapter~\ref{ch:discussion} takes that up;
this chapter's job is the measurement, and the measurement is that the mechanism
is the one the derivation names.

\section{Compute expenditure}
\label{sec:budget}

Every number in this chapter was produced on free cloud GPUs, within a weekly
allowance that the project exhausted more than once. The ladder was designed around that constraint. The accounting of which rung
cost what, which session died and why, and what the constraint forbade outright
is set out in Appendix~\ref{app:harness}. The short version is that compute was the binding constraint throughout.

\part{Assessment}
\chapter{Discussion, limitations, and what would falsify this}
\label{ch:discussion}

\section{What would falsify the thesis}

\begin{table}[htbp]
  \centering
  \caption{The gates, and the response to each. Written before the runs.}
  \label{tab:gates}
  \small
  \begin{tabular}{lp{0.34\linewidth}p{0.44\linewidth}}
    \toprule
    gate & observation that would refute & response \\
    \midrule
    G3 & I-1 shows the Nyquist floor exceeds the estimation term on nearly every primitive in every scene & by Proposition~\ref{prop:reduction} the method is the baseline everywhere. Narrow the claim to mixed-focal and one-sided regimes, or pivot to B2, which this result does not touch. Cost of learning it in week 2: two weeks. In March: the project. \\
    G4 & I-2 shows $\lambda_3/\lambda_1 \approx 1$ on every protocol, including the narrow cone & Equation~\eqref{eq:anisotropy} would be contradicted by measurement, which would indicate an error in the visibility weighting $w_{nk}$ rather than in the geometry. Diagnose before proceeding. \\
    --- & Table~\ref{tab:stress}'s full-orbit row shows a loss against Mip-Splatting & a bug, not a result: Proposition~\ref{prop:reduction} forbids it. The filter never smooths less than the baseline in any direction. \\
    --- & Table~\ref{tab:stress} shows uniform gains across all protocols & also suspicious. The derivation predicts a regime structure; uniform improvement would suggest the comparison is not controlled. \\
    --- & prior work has already made the 3D filter anisotropic or per-direction & the only finding that invalidates the project outright, and the reason the literature due-diligence is scheduled in week 1--2, when it costs nothing, rather than in March. \\
    \bottomrule
  \end{tabular}
\end{table}

\section{Limitations, stated rather than discovered}

\paragraph{First-order approximation.} $\Lambda_k$ is evaluated at the primitive
centre using the affine approximation of the projective map --- the same
linearisation 3DGS makes for the EWA transform. It is consistent with the
renderer; a second-order correction is an ablation, not a fix.

\paragraph{A point-feature likelihood for a volumetric primitive.}
Equation~\eqref{eq:measurement} models the localisation of a point. Real
primitives are extended and semi-transparent, so $\Lambda$ is a proxy for the
photometric information rather than a derivation from it. This is expected to
degrade in textureless regions, where the point-feature model is worst, and that
is a named limitation with a figure.

\paragraph{Pose noise.} $\Lambda$ inherits 3DGS's complete trust in SfM poses.
Because it aggregates over many views, graceful degradation is expected; the
pose-perturbation experiment tests it rather than assuming it.

\paragraph{Compute.} 16\,GB against a stated requirement of 24\,GB removes the
five Mip-NeRF 360 outdoor scenes from the L1 claim. Measured peak VRAM on the
Blender scenes is $\rZeroPeakVram$\,MB, so the constraint does not bind there ---
it binds on the real scenes, and the scene count shrinks accordingly. The three
arguments that carry the thesis --- the exact reduction, the stress suite and the
floor diagnostic --- are all resolution-independent, so the argument survives.

\paragraph{$\beta$ is a prior, not an estimate.} Where the capture constrains a
direction not at all, no amount of estimation theory supplies a scale; $\beta$
imports one. If $\beta$ binds widely, the honest conclusion is that the
estimation limit exceeds what is useful there, and the claim narrows accordingly.

\section{Anticipated questions}

\begin{table}[htbp]
  \centering
  \caption{Questions a viva will ask, and the answers.}
  \small
  \begin{tabular}{p{0.36\linewidth}p{0.56\linewidth}}
    \toprule
    question & answer \\
    \midrule
    Is this Mip-Splatting with more parameters? & It adds two and removes none, and it replaces a scalar chosen from one view with a matrix computed from all of them. By Proposition~\ref{prop:reduction} it reduces to Mip-Splatting exactly when the floor dominates --- a unit test in the repository, not a claim. \\
    Why Fisher information rather than ``resolution''? & Per-view precisions add, and Fisher information is the object that adds. It also carries the right rank structure: one view contributes rank 2, so depth resolution appears only once there is baseline. \\
    You are smoothing --- does that not lose detail? & Never below the floor, which is the baseline's own limit, so lateral detail is untouched by construction. Extra smoothing appears only along unconstrained directions. Training PSNR is expected to fall slightly and held-out PSNR to rise; both are reported. \\
    More views make the filter smaller. Is the anti-aliasing limit not set by the sharpest view alone? & Two filters are in play. The render-side anti-alias filter is set by the sharpest sampling rate --- that is Mip-Splatting's 2D Mip filter, unchanged here. The representation-side limit is an estimation question and genuinely improves with parallax. Conflating them is the error in the scalar 3D filter; the floor keeps them separate. \\
    Did you slow rendering down? & No. The filter changes covariance values, not the rasterisation algorithm, and $\Lambda$ is recomputed ${\sim}290$ times across training, never at render time. Measured fps is in Table~\ref{tab:cost}. \\
    Where does it fail? & Captures so degenerate that $\beta$ binds everywhere, where the filter degenerates to a scene-scale prior; and textureless regions, where the point-feature likelihood poorly models the photometric one. Both are in the limitations section with figures. \\
    \bottomrule
  \end{tabular}
\end{table}

\section{Relationship to the prior course project}
The EE60030 course project (April 2026) argued that 3DGS is sharper than NeRF
because it escapes the MLP's spectral bias. That is correct and incomplete: 3DGS
removed the filter and put nothing in its place. This thesis supplies what should
have replaced it. One correction to that report carries forward --- its Table 3
values are 3DGS Figure 1, a single scene, not Table 1 dataset averages; the
``${\sim}30$\,s/frame'' NeRF figure in it is correct and verbatim from the NeRF
paper.

\input{ch09-schedule}

\appendix
\part{Technical Record}
\input{ch10-derivations}
\chapter{The Reproduction Harness}
\label{app:harness}

\keybox{what this appendix is for}{%
Nothing here is a result. This is the machinery that produced the results: what
runs where, what is checked before anything is measured, how a measurement gets
into the record and how it gets out again, and what the compute allowance
allowed. A reader who takes Chapter~\ref{ch:results} at face value can skip it.
A reader who wants to know whether the numbers can be trusted --- or who wants
to re-run them --- will find the answer here.}

\section{Where The Work Runs}

There is no local GPU. Every training run in this thesis was executed on Kaggle's
free tier: two NVIDIA T4s, $15\,360$\,MiB each, in sessions capped at twelve
hours, against a weekly allowance of $30$ GPU-hours per account. That shaped the
design more than any other single fact, and it is stated first because several
decisions later in this appendix look arbitrary until you know it.

A \emph{rung} is one such session. It clones the repository at a named commit,
builds the CUDA rasteriser from source, converts the dataset, trains, renders,
evaluates, writes one CSV of results and one JSON summary, and prunes everything
else so that the session's output fits in what the platform will hand back. The
scripts are \texttt{kaggle/*.py}; the launcher that pushes them and collects
their output is \texttt{tools/kaggle\_push.py}.

\section{Pre-Flight Assertions}
\label{app:preflight}

An instrument has to be checked before it is used. Twelve assertions run before
the first training step and abort the session if any fails. Each one exists
because the thing it checks can fail silently, and several of them have.

\input{generated/table-preflight}

Three of them:

\begin{itemize}[leftmargin=1.4em,itemsep=3pt]
  \item \textbf{The accelerator check.} Kaggle allocates whichever GPU is free.
        3DGS needs compute capability $7.0$ and a P100 is $6.0$, so a session
        landing on one would fail forty minutes in with a message that looks like
        a code fault. It now aborts in thirty seconds with a marker the launcher
        recognises and relaunches on.
  \item \textbf{The arm-separation diff check.} The 3DGS arm must differ from the
  Mip-Splatting arm by exactly the intended change and nothing else. The check enumerates the files that
        change and refuses anything outside the list. It has fired twice in
        anger --- once on a runtime state file that a careless \texttt{git add}
        had committed, and once when the switch of Appendix~\ref{app:impl}
        reached further than its author believed.
  \item \textbf{The LPIPS backbone.} LPIPS defaults to AlexNet; both papers
        report VGG. The two differ by more than the effects this thesis
        measures. The backbone is pinned, asserted, and written into every row
        of the record as a column, so a row cannot be compared against a
        published figure without the comparison being checkable.
\end{itemize}

\section{The Record}
\label{app:record}

Every measurement lands in one file, \texttt{results/runs.csv}, one row per
(arm, scene, test scale). The file is \textbf{append-only}. A configuration that
is re-measured does not overwrite its predecessor; the predecessor stays and is
marked, in its own \texttt{notes} column, with the name of the stage that
replaced it. Nothing is ever deleted.

That discipline exists because a deleted row is a measurement that cannot be
re-examined, and this project has needed to re-examine several. It also has a
failure mode, which Appendix~\ref{app:defects} describes: a row that is kept but
not marked is a row that is silently \emph{used}. The marking is the only thing
standing between an honest archive and a pooled average, and the marking is now checked --- the generator refuses to emit any table whose
cells draw on more than one build.

\paragraph{From The Record To The Document.} No measured number in this thesis is
typed by hand. \texttt{tools/make\_tex.py} reads \texttt{results/runs.csv} and
writes every table and every inline figure as a LaTeX macro; an unmeasured quantity expands to a visible marker that a reader cannot mistake
for a result. The same loader feeds the figures and the slide deck, so the three cannot
disagree --- and a test asserts that they do not, because for a while they did.

\section{The Apparatus}
\label{app:artefacts}

\begin{table}[htbp]
  \centering
  \caption{What exists, and what each artefact is for.}
  \small
  \begin{tabularx}{\linewidth}{L{0.70}L{1.30}}
    \toprule
    artefact & purpose \\
    \midrule
    \texttt{results/runs.csv} & Append-only. One row per (arm, scene, test-scale, seed), carrying both repository commits, LPIPS backbone, resolution flag, peak VRAM, primitive count, wall-clock and the claimed acceptance level. Every table and figure regenerates from it. \\
    \texttt{tools/make\_tex.py} & Writes every measured table and every inline number in this document. A table with no measured rows is emitted as an explicit placeholder carrying the published targets and nothing in the measured columns; it is never quietly filled from the published column. \\
    \texttt{tools/make\_figures.py} & Generates every figure. In \texttt{--mode measured} it redraws from the CSV with the published curve retained as a faint reference, and refuses to run if the CSV has no usable rows. \\
    \texttt{tools/split\_by\_scale.py} & Recovers the per-scale columns that the shipped evaluation pipeline blends into a single average. Without it there is no multi-scale table. \\
    \texttt{tools/selftest\_pipeline.py} & Runs the whole reporting chain against a fixture with planted, known-in-advance answers, on a machine with no GPU. It is what caught the two silent defects described below. \\
    \texttt{kaggle/*.py} & The rung kernels, each printing every assertion it makes and ending in a single machine-readable summary line. \\
    \bottomrule
  \end{tabularx}
\end{table}

\paragraph{Two Defects The Self-Test Caught.} The
figure generator's ``measured'' mode was drawing the \emph{published} curve under
a measured filename, because \texttt{runs.csv} stores the method as
\texttt{3dgs} while the figures key off \texttt{3DGS}: every lookup missed, the
non-empty guard still passed, and the output was a plausible-looking lie. Second,
the per-scale splitter assumed a JSON nesting that \texttt{metrics.py} does not
write, which would have failed only after both arms had trained. Neither would
have raised an error where it happened. They are recorded here because a
reproducibility claim that has never been tested against a fixture with known
answers is an assertion, not a property.

\section{Cost Per Rung}
\label{app:ladder}

\input{generated/table-budget}

Three entries need a sentence each, because they carry the project's real costs.

\begin{itemize}[leftmargin=1.4em,itemsep=3pt]
  \item \texttt{r3\_real} is the largest single session and did not complete. It
        ran the four indoor Mip-NeRF 360 scenes and was still going when the
        twelve-hour cap arrived. What it produced is kept and reported as a
        four-scene subset, clearly labelled; what it did not produce is not
        estimated.
  \item The 3DGS-arm re-runs after the opacity-compensation fix --- \texttt{c2\_armb} and its continuations
        --- are pure re-measurement cost, incurred because the arm had been built
        wrong. They are the price of the defect in Appendix~\ref{app:defects},
        and they are itemised so that the price is visible.
  \item The stress suite is cheap by comparison. Restricting the training cameras
        does not make training longer, and the sixth protocol --- the one that
        produced the only positive result in this thesis --- trains on three
        cameras and is the fastest run in the table.
\end{itemize}

\paragraph{The Allowance Was Exhausted, Twice.} The first time stalled the
programme for three days and is the reason a second account exists. The binding
constraint on this project has been compute, not code, throughout; the ladder is
ordered so that the cheapest falsifications come first, and the two that could
have ended the project --- does the Nyquist floor bind everywhere, and is the
capture conditioning really anisotropic --- were run before a line of method code
was written.

\chapter{Implementation}
\label{app:impl}

\keybox{what this appendix is for}{%
The thesis compares two published methods and two proposed ones. All four run
from a single checkout, differing only by command-line flags, and this appendix
says exactly what each flag does and where in the code it acts. It is the answer
to ``how do I know the 3DGS arm is really 3DGS?'', and it is a code-level answer.}

\section{Four configurations}

Running each method from its own repository would mean the difference between two
numbers contained the difference between two dataloaders, two metric
implementations and two LPIPS backbones. Every comparison in this thesis is
therefore between configurations of one checkout, and the comparison is the
flags.

\begin{center}\small
\begin{tabular}{llp{0.44\linewidth}}
  \toprule
  configuration & flags & what it is \\
  \midrule
  the Mip-Splatting arm & \texttt{--kernel\_size 0.1} &
    Mip-Splatting as published: 3D smoothing filter on, 2D dilation with the
    opacity compensation. \\
  the 3DGS arm & \texttt{--kernel\_size 0.3 --disable\_3D\_filter
                  --disable\_2D\_mip\_compensation} &
    3DGS's band-limit semantics exactly: no world-space filter, a fixed $0.3$
    screen-space dilation, no opacity compensation. \\
  the anisotropic filter & \texttt{--kernel\_size 0.1 --use\_fisher\_filter} &
    Method I. Replaces the scalar 3D filter with the anisotropic one, floored at
    Mip-Splatting's own value. \\
  the angular mask & post-hoc mask, no training flag &
    Method II. Reads a trained model and the training view directions; masks
    spherical-harmonic degrees. \\
  \bottomrule
\end{tabular}
\end{center}

\section{Two switches for the 3DGS arm}

Mip-Splatting has \emph{two} band-limits and they live in different places.

The first is the 3D smoothing filter: a per-primitive scalar variance added in
world space, computed in Python from the training cameras. Zeroing it is a
one-line change, and with \texttt{--kernel\_size 0.3} --- 3DGS's fixed screen-space
dilation --- the model's world-space semantics are 3DGS's.

The second is the screen-space dilation's \emph{opacity compensation},
$\rho = \sqrt{\det \Sigma' / \det(\Sigma' + kI)}$, which lives in CUDA. Vanilla
3DGS dilates the projected covariance by a fixed $0.3$ and applies no such term.
A Python attribute cannot reach it.

For a long time the 3DGS arm carried only the first switch. That is the defect of
Appendix~\ref{app:defects}: the arm labelled 3DGS was still anti-aliasing with
Mip-Splatting's screen-space term, at exactly the scales where the falsification
gate was scored.

\paragraph{What \texttt{--disable\_2D\_mip\_compensation} touches.} The flag is
threaded from \texttt{arguments/\_\_init\_\_.py} through
\texttt{gaussian\_renderer} into the rasteriser's settings tuple, and read in
both the forward and the backward CUDA kernels:

\begin{itemize}[leftmargin=1.4em,itemsep=2pt]
  \item \texttt{forward.cu}, \texttt{computeCov2D}: $\rho$ is computed only when
        the flag is off; otherwise the coefficient is $1$.
  \item \texttt{backward.cu}, \texttt{computeCov2DCUDA}: the same, and the
        gradient path from opacity through the 2D covariance is removed with it.
\end{itemize}

Because the switch reaches the binary, the pre-flight asserts that the
\emph{compiled} rasteriser carries the field --- checking the source would prove
nothing about what was built.

\section{The degeneracy guard}

Both kernels carry a numerical guard: a primitive whose projected covariance is
singular to the tolerance the determinants are clamped at contributes nothing.
That guard originally sat \emph{inside} the compensation branch, so switching the
compensation off removed the guard as well --- which is not what the switch was
supposed to isolate. It is now hoisted out, in both kernels, and fires
identically in all four configurations.

This is the second lesson of Appendix~\ref{app:defects} and the more general one:
a switch meant to isolate a single term has to be checked against what it leaves
alone, not only against what it changes.

\section{The anisotropic filter in code}

the anisotropic filter needed no CUDA change. The rasteriser already accepts a precomputed
covariance, so replacing a scalar variance with a matrix is a change to how the
covariance is built and not to how it is splatted. Three things happen at
construction, once, before the training loop:

\begin{enumerate}[leftmargin=1.4em,itemsep=2pt]
  \item For each primitive, accumulate
        $\Lambda_k = \sum_n w_{nk}(J_nW_n)^{\!\top}\Sigma_{\mathrm{pix}}^{-1}(J_nW_n)$
        over the training cameras that saw it --- the same Jacobian the renderer
        already forms for the EWA transform.
  \item Invert the spectrum to a covariance floor, and floor \emph{that} at
        Mip-Splatting's own scalar value rather than at a re-derivation of it.
        (Taking theirs is what makes Proposition~\ref{prop:reduction} exact in the
        floor as well as in the covariance; re-deriving it differed by $0.5\,\%$
        at the median on the first call, which is enough to seed a different
        densification trajectory.)
  \item Add it to the primitive's own covariance.
\end{enumerate}

$\Lambda_k$ is never touched inside the loop, which is why rendering is
unaffected (\S\ref{sec:fps}) and why the cost is a construction cost
(\S\ref{sec:methodcost}).

\section{Capture protocols}

A protocol is a rule for selecting a subset of the training cameras. The test set
is never touched --- only the capture the model is fitted from. This is what makes
the stress suite cheap: no new images are rendered and no dataset is regenerated.

The subsets are computed in \texttt{tools/camera\_protocols.py} and passed to the
trainer as an index list. \texttt{mixed focal} additionally downsamples the
selected images, which is a different operation and was, for a while,
not being applied at all --- the protocol passed its index list and nothing else,
so it trained a full orbit under another name. \texttt{tools/test\_camera\_subset.py}
now asserts that each protocol selects what it claims to.

\chapter{Measurement Defects}
\label{app:defects}

\keybox{what this appendix is for}{%
Eight defects were found in this project's measurement, by checks built for the
purpose, and each one produced numbers that looked entirely reasonable while
being wrong. They live in one chapter, away from Chapter~\ref{ch:results}, because a reader
following the argument wants the argument, and a reader auditing the work wants
all the forensics in one place.

Each entry follows the same shape: what the number said, what was actually
happening, how it was found, and what now prevents it. The last of those is the
point. A defect that is fixed but not covered by a check is a defect that will
recur under a different name --- and two of the eight below are the same mistake
at different denominators.}

\section*{The Eight, In One Table}

\begin{center}\small
\begin{tabularx}{\linewidth}{L{1}L{1}L{1}}
  \toprule
  what looked right & what was wrong & what catches it now \\
  \midrule
  The 3DGS arm reproduced 3DGS's published full-resolution figure &
  It still carried Mip-Splatting's screen-space opacity compensation, so it was
  anti-aliasing better than 3DGS at reduced scales &
  The switch is asserted on the \emph{compiled} rasteriser, not the source \\
  \addlinespace
  The opacity-compensation fix switch isolated one term &
  It removed three: the compensation, and a degeneracy guard in each of the
  forward and backward kernels &
  The guard is hoisted out of the branch in both kernels \\
  \addlinespace
  Conditioning and rendered quality agreed per protocol &
  The instrument and the trainer were cloning different branches, so the same
  protocol name meant two different captures &
  Both pin a branch; \texttt{test\_camera\_subset.py} asserts what each protocol
  selects \\
  \addlinespace
  The estimation term exceeded the floor everywhere &
  The constant $0.2$ was carried on one side of the comparison and not the other,
  inflating it fivefold &
  Proposition~\ref{prop:constant} is a unit test, passing to $1.6\times10^{-16}$ \\
  \addlinespace
  Table~\ref{tab:mtmt}'s 3DGS column &
  Every one of its $32$ cells was the mean of two different rasterisers &
  The generator refuses any table whose cells span more than one build \\
  \addlinespace
  Mip-Splatting reproduced to $0.50$\,dB --- L2 &
  A seven-scene mean compared against an eight-scene published figure. The true
  figure is $\stmtArmAMaxDelta$\,dB --- L1 &
  Against-published uses the arm's own scene set; between-arms uses the matched
  set; a test asserts all three outputs agree \\
  \addlinespace
  The offset was explained by GOF densification &
  Asserted for months, never tested. Measured, the effect is
  $\gofMaxEffect$\,dB and not consistently signed &
  The pre-GOF commit is a probe in the record (Table~\ref{tab:pregof}) \\
  \addlinespace
  The angular mask lost heavily to magnitude pruning &
  Against a \emph{free} control with more freedom than the angular mask has. At matched
  structure the margin is $\btwoBlockedGapMin$--$\btwoBlockedGapMax$\,dB &
  The control carries the angular mask's own nested-degree structure \\
  \bottomrule
\end{tabularx}
\end{center}

\medskip
Five of the eight are given in full below. The other three --- the denominator,
the densification explanation and the angular mask's unfair control --- cannot be separated
from the results they correct, so they are told where those results are:
\S\ref{sec:table1}, \S\ref{sec:densification} and \S\ref{sec:btwo}
respectively.

\paragraph{One Pattern, Six Times.} Six of the eight are the same error: two
quantities were compared that were not measured over the same thing. Two
rasterisers averaged into one cell. A seven-scene mean against an eight-scene
figure. A protocol name meaning two captures. A constant present on one side and
not the other. A pruning control with more freedom than the method it judged. In
every case the arithmetic was right and the comparison was not, and in every case
the result looked plausible --- which is why none of them was caught by reading
the numbers, and all of them were caught by a check that asked what the numbers
were means \emph{over}.

\section{The Fourth Check: The Rasteriser On The 3DGS Arm}

The single-codebase construction of \S\ref{sec:onecodebase} gives both arms
Mip-Splatting's rasteriser, and sets the 3DGS arm apart by zeroing \texttt{filter\_3D}
and taking \texttt{kernel\_size}~$=0.3$. That reproduces 3DGS's \emph{3D}
band-limit exactly --- which is what check~2 verifies --- and says nothing about
the \emph{2D} one. Mip-Splatting's screen-space filter dilates the projected
covariance by $k$ and then compensates the opacity by
\begin{equation}
  \rho \;=\; \sqrt{\frac{\lvert\Sigma'\rvert}{\lvert\Sigma' + kI\rvert}} ,
\end{equation}
so that dilation does not brighten the primitive. Vanilla 3DGS applies the same
fixed $0.3$ dilation and \textbf{no compensation at all}. Setting
\texttt{kernel\_size}~$=0.3$ reproduces the dilation magnitude and leaves $\rho$
in place, so the 3DGS arm kept an anti-aliasing term that 3DGS does not have --- and
the term it kept is precisely the one that protects a dilated primitive from
losing energy as the test sampling rate falls. An arm that anti-aliases better
than 3DGS degrades less than 3DGS, which is what the scale-collapse gate measured.

\paragraph{The Test That Separates The Two Explanations.} $\rho$ is
computed in CUDA, so the switch is a rasteriser change: a
\texttt{mip\_compensation} flag threaded through
\texttt{Gaussian\allowbreak Rasterization\allowbreak Settings} into
\texttt{computeCov2D}, which skips
the $\rho$ term while keeping the dilation. Everything else in that function is
identical to \texttt{graphdeco-inria/\allowbreak diff-gaussian-\allowbreak rasterization} line
for line.
The build asserts the compiled binary carries the field before any training
starts, because checking the source would prove nothing about the binary.

One scene, \texttt{lego}, 30\,000 iterations, everything else held fixed:

\input{generated/table-c1}

Full resolution barely moves: $\cOneBaseFull$ against $\cOneFull$, which is
what must happen, because at the training sampling rate the dilation is
negligible against the primitive's own extent and $\rho \approx 1$. Every
reduced scale collapses, and lands on the shape of the published 3DGS row:
$\cOneHalf$ against $26.95$, $\cOneQuarter$ against $21.38$, $\cOneEighth$
against $17.69$. The compensation was the whole of the defect.

\keybox{the densification hypothesis}{%
Both columns of Table~\ref{tab:c1} ran GOF's densification. If densification were
the cause of the shortfall, the right-hand column could not have collapsed --- and
it collapses by $\cOneCollapse$\,dB at $\tfrac18$, landing within $0.3$\,dB of
the published figure at every reduced scale. The earlier diagnosis is therefore
wrong, and it is left standing in this chapter because how a wrong diagnosis
survived matters: it was never given an experiment that could refute
it. It explained the anomaly as a property of the codebase, which is the most
comfortable kind of explanation and the one that most deserves a test.

One point is easy to read the wrong way. The comparison that settles this runs
\emph{between the two columns}; either column set against the published row
settles nothing. Both columns are \texttt{lego} and
the published row is an eight-scene mean, so the agreement at reduced scales is
corroboration and not a measurement. The measurement is that switching one CUDA
term moved $\tfrac18$ by $\cOneCollapse$\,dB and full resolution by less than
$0.05$.}


\section{What The Vanilla Arm Cost In Robustness}
\label{sec:shipfails}

Removing the compensation also removed a guard. The same CUDA branch that
computes $\rho$ carries \texttt{coef = 0} for a primitive whose projected
covariance is singular to the tolerance the determinants are clamped at --- a
numerical guard, not part of the anti-aliasing model. With the compensation
switched off, degenerate primitives survived at full opacity with an unbounded
screen-space radius, and two scenes died reproducibly, on both GPUs, with a CUDA
illegal memory access around iteration $3\,000$. Six trained normally, which is
what made it look transient.

Hoisting the guard out of the branch fixed one of the two. The other then failed
in the \emph{backward} pass instead, because the backward kernel still had its
copy of the guard inside the branch: the forward was contributing nothing for a
degenerate primitive while a gradient still flowed to its opacity. Zeroing a
forward contribution and leaving its gradient live is not a consistent pair, and
matching them is what the opacity-compensation fix switch should have done from the start --- it was
supposed to isolate one term and it removed three.

\paragraph{\texttt{ship} Does Not Complete, And Is Reported As Such.} With both
guards matched it fails a third time, and a fourth under instrumentation. The
first three tracebacks all pointed at \texttt{ssim}, in the \emph{next}
iteration's loss, which is a red herring: a CUDA fault is reported
asynchronously, at whichever call next touches the device. Re-run with
\texttt{CUDA\_LAUNCH\_BLOCKING=1} so that each launch is synchronous, it
localises instead to the forward rasteriser --- the fault surfaces on the first
tensor derived from its output, \texttt{radii > 0} in
\texttt{gaussian\_renderer}, with the backward pass never reached.

That clears three explanations. The backward kernel, which the second failure
had suggested, is clear: the fault surfaces before the backward pass is reached.
Memory exhaustion is clear: every attempt raises an illegal address and none
raises an allocation error, and \texttt{ship} completes on the Mip-Splatting arm
on the same $15\,360$\,MiB card at a measured peak of $13\,755$\,MiB. And the
onset moves between attempts, which clears the deterministic fault at a fixed
point this section claimed before it was instrumented.

The scene has now been attempted nine times on the vanilla arm --- four
diagnostic runs and then five consecutive attempts in one session, since a fault
whose onset moves is a fresh draw each time, so retrying it is a measurement. \textbf{All nine failed}, every one with
\texttt{cudaErrorIllegalAddress} and none with an allocation error, at
iterations
\[
  3\,000,\quad 7\,220,\quad 8\,580,\quad 10\,100,\quad 10\,210,\quad
  10\,320,\quad 11\,820,\quad 12\,100
\]
and one earlier still. What varies is \emph{when}, not \emph{whether}: the
failure is certain and its timing is not. That settles the question retrying was
asked to settle --- there is no lucky run to be had --- and it is why the scene
is reported as failed and left there.

What remains is an out-of-bounds access in the forward raster whose onset
depends on the optimisation trajectory, on much the largest model in the suite:
\texttt{ship} reaches $442\,235$ primitives on the Mip-Splatting arm, more than any other
Blender scene, at $\sim\!90\,\%$ of the card.

It is excluded from the single-scale 3DGS-arm rows, and because every
between-arm
table in this chapter is restricted to the scenes both arms have, it is excluded
from the Mip-Splatting arm's there too --- a seven-scene mean set against an eight-scene one compares two different things. Nothing was reduced to make it fit: the protocol's rule for a scene
that will not run is to record it and move on, and that rule does not become
optional when the scene is one of eight, as it would be one of nine.

Two things keep the exclusion narrow. \texttt{ship}
trains without incident on the \emph{multi-scale} protocol, on all four scales,
so Table~\ref{tab:mtmt} is a full eight-scene comparison and the failure is
specific to single-scale training. And it trains without incident on the Mip-Splatting arm, so
the fragility belongs to the vanilla band-limit semantics, with the scene, the
data and the harness all cleared. Both point the same way: removing the
compensation costs robustness on the one scene that sits closest to the limits of
the card, and it does so through the forward raster.

The mechanism is \emph{not} settled beyond that, and the earlier draft of this
section overstated it. Hoisting the degeneracy guard out of the compensation
branch was a real defect and a real fix --- it is what recovered \texttt{chair}
--- but it cannot also be the explanation for the failures that persist after
it, because the guard now fires identically whether the compensation is on or
off. What can be said from measurement is the conjunction above: forward raster,
varying onset, no allocation error, largest model in the suite, vanilla arm
only. Naming a cause more specific than that would be inference presented as
result, which is the thing \S\ref{sec:g2} already cost this project once.

\keybox{what makes a switch a control}{%
Each of those three removals was invisible: no assertion covered them, the
numbers that came back were plausible, and six of eight scenes were unaffected.
What exposed them was a failure that could not be waved through --- a crash --- and
then a failure that \emph{moved} from the forward pass to the backward one when
half the fix was applied, which is what identified the second half. A switch
meant to isolate a single term has to be checked against what it leaves alone,
not only against what it changes.}


\section{A Column Averaged Over Two Rasterisers}
\label{sec:blend}

The 3DGS arm's column here was re-measured after the opacity-compensation fix, and for a while afterwards each of
its $32$ cells averaged the old measurement together with the new one, where the
new one alone belonged there. The record is append-only by design (Table~\ref{tab:levels}): a
re-measured configuration leaves its predecessor in the file, marked, and the
tables skip anything marked. The marking rule identified the rows to retire by
the rung name in their notes --- and the re-run was the same rung, so its rows
carried the same name, the rule exempted both, and nothing was retired. Both
rows then satisfied the table's selector and were averaged.

Nothing about the output looked wrong. The two builds differ only in whether the
degeneracy guards in the forward and backward passes match, which moves a scene
by hundredths of a decibel on a capture this well conditioned, so the blended
column sat exactly where a plausible column sits. The same defect had quietly
taken hold in the anisotropic filter parity rows, where the two builds differed in whether the
filter floor was Mip-Splatting's own value or a re-derivation of it.

What was missing was not a better rule but any check at all on the invariant the
rule exists to maintain. There is one now: the generator refuses to emit a table
if any reported cell draws on more than one build, and it names the cell. It is
four lines, it would have caught both instances the day they appeared, and it
subsumes every future variant of the same mistake --- which a third special case
in the marking rule would not have.

\keybox{what append-only applies to}{%
Keeping superseded rows is what makes a result auditable, and this project has
been strict about it. But a row that is kept and not marked is a row that is
\emph{used}, and the marking is the only thing standing between an honest
archive and a silently pooled average. A discipline that depends on a marking
step needs the pooled average checked, not the marking step trusted.}


\section{Two Kernels, One Set Of Protocol Names}
\label{sec:branchdefect}

The instrument results and the PSNR results of \S\ref{sec:table3} were, for a
time, measurements of \emph{different captures} reported under the same protocol
names, and the mechanism matters, because nothing in either run could have
detected it.

The instrument kernel cloned the repository's default branch. The training kernel
clones the method branch. Between them, the definition of \texttt{arc} and
\texttt{cone} had changed: the default branch sizes them by angular \emph{width},
giving $20$ cameras across $81^\circ$ and $3$ across $30^\circ$; the method
branch sizes them by camera \emph{count}, giving $25$ across $93^\circ$ and $10$
across $54^\circ$. The change was deliberate and its stated reason is sound --- a
$20^\circ$ wedge catches three cameras on a Blender capture, and three views turns a conditioning experiment into a sparse-reconstruction
one. What was missed
is that only one of the two kernels was updated.

Each run was internally consistent, each printed the span it had built, and each
was correct about its own capture. The error was in reading them side by side,
and it survived because both numbers looked reasonable: a cone \emph{is} narrower
than an arc under either definition, so the ordering the chapter argues for held
either way.

\keybox{the cost of this defect}{%
It cost the first reading of the floor-binding and anisotropy gates --- and it bought the \texttt{pencil}
protocol. Rather than choose between the two definitions, both are kept: the
count-based \texttt{cone} at $51^\circ$, which is a trainable subset, and the
width-based \texttt{pencil} at $20^\circ$, which is the regime the derivation is
about. The crossover between them is where the estimation term overtakes the
floor, and with both in the table that crossover is measured directly. The instrument kernel is now pinned to the branch the training
kernel clones, and every protocol table carries the span the session that
produced the row actually built.}

\section{A Constant That Was Carried On One Side Only}
\label{sec:m2defect}

One more defect has to be stated before the numbers, because it changed them.
The estimation term is $(s\,\sigma_{\mathrm{pix}})^2/\lambda_i$, and
Proposition~\ref{prop:constant} fixes $(s\,\sigma_{\mathrm{pix}})^2 = 0.2$ ---
the same constant Mip-Splatting's floor $f_k^2 = 0.2\,(d_{\min}/f_{\max})^2$ is
built from. Both sides of the floor-occupancy instrument comparison therefore carry it. The instrument
carried it on the floor and dropped it from the estimation term, which inflated
that side by a factor of five, or $\sqrt{5} = 2.24$ in $\sigma$. The filter's own
implementation had always applied it.

The symptom was two measurements of one quantity disagreeing --- on the arc and
the cone, and not at the extremes, where a factor of five changes no verdict.
Nothing raised, and both numbers were plausible on their own.

\keybox{Proposition~1 as a test}{%
The identity that fixes the constant is checkable: at one view, on the optical
axis, the estimation term must \emph{equal} the floor. It now does, to
$1.6\times10^{-16}$, and \texttt{tools/test\_instruments.py} asserts it on every
run. Without the $0.2$ it is exactly five times the floor, so the test cannot
pass either way by accident. This is the over-determination the audit's M-2
named --- ``$s$ defaults to 1'' and ``$s\,\sigma_{\mathrm{pix}} = \sqrt{0.2}$''
cannot both be free --- surfacing as a live numerical error, the most useful form it could take; as an
ambiguity in the prose it would have stood for years.}


\section{Three Bugs The Parity Requirement Caught}
\label{sec:b1bugs}

The first implementation of the anisotropic filter reported exactly what
\S\ref{sec:instrumentresults} predicts internally --- the floor binding on
$100\,\%$ of primitives, the filter isotropic to two decimals --- and still
differed from Mip-Splatting by up to $1.22$\,dB, thirty times the seed spread of
\S\ref{sec:spread}. Section~\ref{ch:discussion} names that case in advance: a
loss on the full-orbit row is a bug, not a result, because
Proposition~\ref{prop:reduction} forbids it. Diagnosing it took three passes, and each step ruled something out.

\begin{enumerate}[leftmargin=1.4em,itemsep=3pt]
  \item \textbf{The floor.} The anisotropic filter re-derives $f_k$ from the camera geometry, so the
        first suspect was that its floor was not Mip-Splatting's. A cross-check
        was added --- \texttt{compute\_3D\_filter} runs alongside and the two are
        compared on every call --- and they agree to $1.16\times10^{-10}$
        absolute, $10^{-7}$ relative, for the whole run. Not the cause. The anisotropic filter now
        takes their floor verbatim anyway, since \S\ref{sec:proposition2} floors
        at ``Mip-Splatting's own value'' and a re-derivation is a needless
        opportunity to differ.
  \item \textbf{The covariance path.} The anisotropic filter passed $\Sigma + \Sigma_{\mathrm{filt}}$
        through the rasteriser's precomputed-covariance input while
        Mip-Splatting passes scales and rotations. The two are algebraically the
        same, but they are different code. Where the filter is isotropic on
        every primitive, the anisotropic filter now takes Mip-Splatting's own path ---
        Proposition~\ref{prop:reduction} used as a computational identity in the code,
        and no longer only as an argument on paper. The gap did not move: $1.245$\,dB. Not the
        cause either.
  \item \textbf{The opacity.} What both paths share is
        Equation~\eqref{eq:filter}'s compensation, and it contained
        \texttt{clamp\_min(1e-12)} on both determinants. $|\Sigma|$ is
        $\prod s_i^2$: a primitive with $s = 0.005$ has $|\Sigma| = 1.6\times
        10^{-14}$, and one with $s=0.001$ has $10^{-18}$. Both determinants
        clamped to the same floor, their ratio became exactly $1$, and the
        compensation switched itself off for most of the model. The anisotropic filter was rendering
        with uncompensated opacity while every baseline number in this thesis
        was produced with it --- which is precisely the observed signature:
        denser and sharper at full resolution, worse at every reduced scale.
\end{enumerate}

With the clamp reduced to a division guard, $|\Sigma|$ taken from the scale
vector where $\prod s_i^2$ is exact, and the isotropic regime using
Mip-Splatting's own scalar expression, the largest difference over the two scenes
then measured fell from $1.245$\,dB to $0.136$\,dB. Over the
$\methodParityScenes$ scenes for which a \emph{paired} measurement exists it is
now $\mathbf{\methodParityMax}$\,dB, comfortably below the seed spread of
\S\ref{sec:spread}.

\keybox{what ``paired'' has to mean here}{%
Proposition~\ref{prop:reduction} says the anisotropic filter \emph{becomes} Mip-Splatting when the
floor dominates, so the two numbers being subtracted have to come from the same
capture, the same harness invocation and the same build; otherwise the
difference also contains whatever separates two runs. That is a stricter
requirement than it sounds, and this table did not meet it. The selector behind
the macro accepted the standard single-scale rows as a stand-in for
Mip-Splatting at the full protocol, so the anisotropic filter was being differenced against a sweep
from three commits earlier that had never used the protocol, with two scenes
entering twice on one side and once on the other. It reported $0.021$\,dB. The
figure was not wrong by much, and that is the point: a comparison can be
mis-specified and still land on a plausible number, which is why the pairing is
now enforced in the generator, and why the scene count is a macro that reads
from the record.}

\keybox{what made this findable}{%
The unit test of \S\ref{sec:proposition2} passed throughout, because its scales
were $0.01$--$0.06$, where $\prod s_i^2$ sits at the clamp and the fault is
nearly invisible. A test whose inputs do not span the real range will certify a
broken component. It now covers $|\Sigma|$ down to $1.4\times10^{-18}$ and prints
that range in its own output. The bug itself was silent --- no error, plausible
numbers, the wrong shape --- and what exposed it was requiring parity to hold \emph{exactly}.}


\end{document}
